\documentclass[11pt,a4paper]{ctexart}

\usepackage[margin=2.35cm]{geometry}
\usepackage{amsmath,amssymb,mathtools,bm}
\usepackage{booktabs,array}
\usepackage{enumitem}
\usepackage[numbers,sort&compress]{natbib}
\usepackage[colorlinks=true,linkcolor=blue,citecolor=blue,urlcolor=blue]{hyperref}
\setlength{\emergencystretch}{2em}

\newcommand{\dd}{\mathrm d}
\newcommand{\ii}{\mathrm i}
\newcommand{\Id}{\mathbb I}
\newcommand{\cV}{\mathcal V}
\newcommand{\cE}{\mathcal E}
\newcommand{\cH}{\mathcal H}
\newcommand{\cM}{\mathcal M}
\newcommand{\cB}{\mathcal B}
\newcommand{\cR}{\mathcal R}
\newcommand{\cN}{\mathcal N}
\newcommand{\Emb}{\operatorname{Emb}}

\title{MadStree v0.11: tree formula}
\author{\texttt{package-MadStree} 公式手册}
\date{2026-08-14}

\begin{document}
\maketitle

\begin{abstract}
本手册对应 \texttt{package-MadStree} v0.11；time-only sector 继续由 root propagator 顺序的定长字符串标识。
本文是 \texttt{package-MadStree} 的公式手册。该程序包面向任意 massive/massless 混合 dS 树图，不生成一般 IBP 方程组，而由 topology 直接组装主积分、迭代约化和 block-triangular dlog 微分方程；具体使用指南在公开接口稳定后补入。本文从局部函数系统的张量积结构重新推导所需公式。Chen--Feng 的纯 massive 公式可统一理解为若干一维无 theta 指数系统与二维 Hankel 系统的 Kronecker 和；每个 $\ii\chi_bq_b$ 都是一维指数生成元与其余单位矩阵的张量积。含 theta 的 massless full propagator 不能等同于无 theta 的单指数因子：约化前，它在两个端点同样具有 $2\otimes2=4$ 个 Hankel 状态；利用 massless 的 $00/11$、$10/01$ 关系后，这四态被一个常数投影压成整条边共享的二维状态。该二维因子同时出现在两个顶点的 time-IBP 中，故按顶点独立的 tensor factorization 失效，但全图上的 tensor-product/Kronecker-sum 构造仍然成立。所有 regular IBP、迭代矩阵和 dlog 对角块因而可由 topology 直接写出；theta 导数只增加连接缩边 sector 的固定非对角块，无需先生成大规模 IBP 系统再反推出公式。
\end{abstract}

\tableofcontents

\section{范围、参考 convention 与结论}

本文以文献 \cite{Chen:2023eic} 的式 (3.1)--(3.5)、(3.37)、(3.47)、(3.50)、(3.54)--(3.55) 和 (3.64)--(3.68) 为公式来源。该文使用 $\tau^{\nu_0}$ 以及
\begin{equation}
 h(\nu,0;z)=z^{-\nu}H_{\nu}(z),
 \qquad
 h(\nu,1;z)=\partial_z h(\nu,0;z),
 \qquad z=-k\tau,
 \label{eq:paper-h-definition}
\end{equation}
以及自然二进制 basis。MadStree 始终规定 $\nu=|\nu|$，只使用一个函数名 $h$，并定义
\begin{equation}
 \boxed{h(\nu,0;z)=z^{s_\nu\nu}H_\nu(z)},
 \qquad s_\nu\in\{+1,-1\}.
 \label{eq:positive-order-h}
\end{equation}
缺省 \texttt{NuConvention -> "Positive"} 对应 $s_\nu=+1$；\texttt{"Negative"} 对应 $s_\nu=-1$，也就是论文 convention。两种定义并非只差常数；它们相差 $z^{2\nu}$。但是 Hankel 方程只含 $\nu^2$，代入后两套无量纲 EOM 分别为
\begin{equation}
 h''+\frac{1-2\nu}{z}h'+h=0\quad(s_\nu=+1),
 \qquad
 h''+\frac{1+2\nu}{z}h'+h=0\quad(s_\nu=-1).
\end{equation}
所以它们的全部公式严格只差 $\nu\to-\nu$。也可以把正 prefactor basis 与 2401 convention 中取负参数的同一个函数 $h(-\nu,0;z)=z^\nu H_{-\nu}(z)$ 比较；$H_{-\nu}$ 与 $H_\nu$ 只差运动学无关的 branch phase。下文引用 2401 的矩阵公式时仍把其中的参数写成 $\nu$；MadStree 的正 prefactor缺省实现统一执行
\begin{equation}
 \boxed{\left.\nu\right|_{2401}\longrightarrow-\nu}.
 \label{eq:paper-to-positive-order}
\end{equation}
这不是只替换 Hankel 下标；EOM、Wronskian、$M_1$、contact 幂次、递推和 dlog 中的每个 $\nu$ 都必须同时替换。最后一节另给出到 package 的 $(-\tau)^A$、normalized-$J$ 和 sector normalization 适配。

必须先区分两类对象：
\begin{itemize}
  \item 外腿或 $G^{+-}/G^{-+}$ 中已经选定方向的纯指数没有 theta，只是一维函数系统；
  \item $G^{++}/G^{--}$ 的 massless full propagator 含两项和 theta。约化前，它与 massive full propagator 一样，在两个端点各有一个二维 $h$ 状态，合计四态。
\end{itemize}
本文的核心结果是：massless 额外关系把第二类对象的四态压成一个共享二维因子，而不是把它误当成两个独立的一维指数。

公式层结论分为两级。固定 sector 的 regular 矩阵、四态到二态的投影以及共同对角化是直接代数结论；contact 对各 lower sector 的注入结构同样可直接写出，但任意 mixed/common-theta normalization 下的 constant-residue dlog 与 flatness 仍须独立验证。因此本文不给现有 package 的 \texttt{PendingRederivation} 自动解锁。

\section{局部函数系统：一维指数与二维 h}
\label{sec:local-function-system}

\subsection{无 theta 指数的一维系统}

不是每个顶点只允许一个固定名字的 $k_0$. 对每个不带 theta 的指数 building block $b$，记录
\begin{equation}
 E_b(\tau_{v(b)})
 =\exp\!\left(\ii\chi_bq_b\tau_{v(b)}\right),
 \qquad \chi_b\in\{+1,-1\},
 \label{eq:phase-exponent-block}
\end{equation}
其中 $q_b$ 保留用户输入的名字和定义。纯顶点能量、massless cross/external 分支都可以是这种 block；同一个顶点可以关联多个 $b$. 把 $E_b$ 视为一维空间 $\cE_b\simeq\mathbb C$ 上的向量，其 time generator 是
\begin{equation}
 Q_b=[\ii\chi_bq_b],
 \qquad I_{\cE_b}=[1].
 \label{eq:exponential-1d}
\end{equation}
因此这些 blocks 不增加主积分维数，但必须分别编号，以便保留用户变量、求导 Jacobian 和边界坐标。对顶点 $v$，矩阵层只需使用 signed sum
\begin{equation}
 p_{0,v}=\sum_{b:\,v(b)=v}\chi_bq_b,
 \qquad
 \sum_{b:\,v(b)=v}\Emb_b(Q_b)=\ii p_{0,v}I.
 \label{eq:vertex-signed-phase-energy}
\end{equation}
论文的 $e^{\ii k_{0,v}\tau_v}$ 是只有一个 block、$\chi_b=+1$、$p_{0,v}=k_{0,v}$ 的特例。每个 Hankel 导数项也都乘着所有一维 blocks 的单位矩阵，所以所有 regular time derivatives 仍是同一种 Kronecker-sum 项。

\subsection{二维 h 系统}

以下先沿用 2401 矩阵公式中的参数 $\nu$。令
\begin{equation}
 \bm h(\nu;z)=
 \begin{pmatrix}h(\nu,0;z)\\h(\nu,1;z)\end{pmatrix},
 \qquad
 P_1=\begin{pmatrix}0&0\\0&1\end{pmatrix},
 \qquad
 \sigma_2=\begin{pmatrix}0&-\ii\\ \ii&0\end{pmatrix}.
\end{equation}
文献式 (3.2)--(3.4) 等价于
\begin{equation}
 \partial_\tau\bm h
 =\left[-\ii k\sigma_2-\frac{2\nu+1}{\tau}P_1\right]\bm h.
 \label{eq:h-first-order-system}
\end{equation}
其中 $-\ii k\sigma_2$ 是与时间无关的 $2\times2$ 矩阵，$-(2\nu+1)P_1/\tau$ 与时间幂次的降低合并进入 $M_1$。

同一系统在无量纲变量 $z=-k\tau$ 中满足
\begin{equation}
 \partial_z^2h(\nu,0;z)
 +\frac{2\nu+1}{z}\partial_zh(\nu,0;z)
 +h(\nu,0;z)=0.
 \label{eq:h-dimensionless-eom}
\end{equation}
这里的常数项是 $1$，不是 $k^2$。$k$ 只由
\begin{equation}
 \partial_\tau=-k\partial_z
 \label{eq:tau-z-chain}
\end{equation}
恢复。这一分层是后文复用论文矩阵的关键。

\subsection{正/负 prefactor convention 与 massless 指数核}

Hankel 方程只依赖阶数平方，但这里固定 Hankel order 为 $\nu=|\nu|$，只切换 $z^{\pm\nu}$。prefactor 和公式参数必须按式 \eqref{eq:paper-to-positive-order} 一起变换。在 MadStree 的缺省正 prefactor convention 中，当 $\nu=1/2$ 时，
\begin{equation}
 h(1/2,0;z)\big|_{s_\nu=+1}=z^{1/2}H_{1/2}^{(1,2)}(z)
 \propto e^{\pm\ii z},
 \qquad
 \left.(2\nu+1)\right|_{2401}\xrightarrow{\nu\to-\nu}1-2\nu=0.
 \label{eq:massless-h-eigenstate}
\end{equation}
于是
\begin{equation}
 \partial_z^2h(1/2,0;z)+h(1/2,0;z)=0\quad(s_\nu=+1),
 \label{eq:massless-z-eom}
\end{equation}
且第一导数态是零阶态的分支本征态，
\begin{equation}
 h^{(1)}(1/2,1)=+\ii h^{(1)}(1/2,0),
 \qquad
 h^{(2)}(1/2,1)=-\ii h^{(2)}(1/2,0).
 \label{eq:h-branch-eigenvalues}
\end{equation}
负 prefactor 选项给出 $z^{-1/2}H_{1/2}^{(1,2)}\propto z^{-1}e^{\pm\ii z}$，不再是纯指数；其公式由正 prefactor 结果作 $\nu\to-\nu$ 得到。纯虚阶情形应先选定平方根分支；这里的 $|\nu|$ 是 convention 标签，不表示在代数公式中调用复数 \texttt{Abs[nu]}。

\section{纯 massive 顶点公式本来就是 Kronecker 和}

考虑论文的 $p$-fold 顶点积分
\begin{equation}
 V_v(\nu_{0,v},\bm a)
 =\int\dd\tau_v\,
 \tau_v^{\nu_{0,v}}e^{\ii p_{0,v}\tau_v}
 \prod_{i=1}^{p}h(\nu_i,a_i;-k_i\tau_v),
 \qquad a_i\in\{0,1\}.
\end{equation}
它的局部空间是
\begin{equation}
 \cV_v
 =\bigotimes_{b:\,v(b)=v}\cE_b
 \otimes\bigotimes_{i=1}^{p}\cH_i,
 \qquad
 \dim\cE_b=1,
 \quad
 \dim\cH_i=2.
 \label{eq:vertex-tensor-space}
\end{equation}

对任一局部因子 $\alpha$，定义嵌入
\begin{equation}
 \Emb_\alpha(A)
 =I_1\otimes\cdots\otimes A_\alpha\otimes\cdots\otimes I_N.
 \label{eq:embedding-definition}
\end{equation}
于是对被积函数做一次 time IBP，论文式 (3.37) 可直接写成
\begin{align}
 M_{1,v}(\nu_{0,v})
 &=\nu_{0,v}I_v
 -\sum_{i=1}^{p}(2\nu_i+1)\Emb_i(P_1),
 \label{eq:massive-M1-projector}\\
 M_{0,v}
 &=\ii p_{0,v}I_v
 +\sum_{i=1}^{p}\Emb_i(-\ii k_i\sigma_2).
 \label{eq:massive-M0-kronecker}
\end{align}
第一式使用 $P_1=(I-\sigma_3)/2$ 后，正好等于论文写法
\begin{equation}
 M_{1,v}
 =\sum_i\left(\nu_i+\frac12\right)\Lambda_3^{(i)}
 +\left(\nu_{0,v}-\frac p2-\sum_i\nu_i\right)I_v.
\end{equation}
因此 $p_0I$ 与 $k_i\Lambda_2^{(i)}$ 没有本质区别：前者是一维指数 generators 的和，后者来自二维 h 因子；二者都是局部 generator 乘所有 spectator identities。

令 $\bm f_v(a)$ 表示把 $\nu_{0,v}$ 平移整数 $a$ 后的 $2^p$ 维积分向量。忽略 theta remaining term 时，
\begin{equation}
 M_{1,v}(\nu_{0,v}+a)\bm f_v(a-1)
 +M_{0,v}\bm f_v(a)=0.
 \label{eq:massive-vertex-ibp}
\end{equation}
降低一步直接为
\begin{equation}
 A_{-,v}(a)
 =-M_{1,v}(\nu_{0,v}+a)^{-1}M_{0,v}.
 \label{eq:Aminus-local}
\end{equation}

论文的常数矩阵
\begin{equation}
 T=\frac1{\sqrt2}
 \begin{pmatrix}1&-\ii\\-\ii&1\end{pmatrix},
 \qquad
 T\sigma_2T^{-1}=\sigma_3,
 \label{eq:paper-T}
\end{equation}
同时对角化所有 $M_0$ 中的二维局部 generator。令 $T_v=T^{\otimes p}$，则
\begin{equation}
 \widetilde M_{0,v}=T_vM_{0,v}T_v^{-1}
\end{equation}
为对角矩阵，升高一步为
\begin{equation}
 A_{+,v}(a)
 =-T_v^{-1}\widetilde M_{0,v}^{-1}T_v
 M_{1,v}(\nu_{0,v}+a+1).
 \label{eq:Aplus-local}
\end{equation}
所以纯 massive 的公式从一开始就不需要一般大矩阵求逆：$M_1$ 在原 basis 中对角，$\widetilde M_0$ 在 $T$ basis 中对角。

\section{从逐顶点空间到全图空间}

设固定 sector $s$ 的 active vertices 为 $V_s$。在尚未使用 massless 额外关系以前，每个传播子端点的 h 状态都是独立二维因子；每个无 theta 指数是一个一维因子。因此 regular 全图空间可写成
\begin{equation}
 \cV_s^{\rm unreduced}
 =\bigotimes_{b\in E_0(s)}\cE_b
 \otimes
 \bigotimes_{\text{all active endpoints }i}\cH_i.
 \label{eq:global-unreduced-space}
\end{equation}
这里 $E_0(s)$ 是 active phase-exponent block 集。每个顶点的 $M_{0,v}$、$M_{1,v}$ 都由式 \eqref{eq:embedding-definition} 嵌入这个共同空间。对纯 massive 图，或对没有 theta 的 cross propagator，顶点 $v$ 的 regular operator 只作用于属于该顶点的因子，所以不同顶点的 vertex families 可以独立因子化。$G^{++}/G^{--}$ 的 theta 导数不属于 regular operator，而作为 lower-sector remaining term 另行加入。

这个全图写法比“先为每个顶点各写一次公式”更基本。后者成立，只是因为纯 massive 情形中不同顶点的 regular operator 支撑在互不相交的二维因子上。

\section{massless full edge：四态到共享二态}

\subsection{它不是一维指数因子}

考虑一条含 theta 的 massless full edge $e=(u_e,v_e)$。在使用额外关系以前，它与 massive propagator 完全一样：端点 $u_e$ 有二态 $a\in\{0,1\}$，端点 $v_e$ 有二态 $b\in\{0,1\}$，所以边的局部空间是
\begin{equation}
 \cH_{e,u}\otimes\cH_{e,v}\simeq\mathbb C^4.
\end{equation}
按论文自然顺序取
\begin{equation}
 \bm F_e^{(4)}
 =(F_{00},F_{01},F_{10},F_{11})^T.
\end{equation}
由式 \eqref{eq:h-branch-eigenvalues}，在论文无量纲 h-state convention 下，
\begin{equation}
 F_{01}=-F_{10},
 \qquad
 F_{11}=F_{00}.
 \label{eq:h-state-massless-relations}
\end{equation}
这里没有遗漏 $k^2$：每个 $F_{ab}$ 使用的是 $h_1=-k^{-1}\partial_\tau h_0$，而不是未除去 $k$ 的物理端点导数。它也不是 package 抽出每次端点导数相位后的代数 quotient；后者的 $11$ 关系带负号。物理关系在第 \ref{sec:physical-k2} 节恢复，package quotient 在第 \ref{sec:package-adapter} 节给出。

选择共享二态
\begin{equation}
 \bm g_e=(F_{00},F_{10})^T\in\cM_e\simeq\mathbb C^2.
\end{equation}
四态与二态之间的常数 embedding 和 projection 为
\begin{equation}
 \boxed{
 \bm F_e^{(4)}=S_e\bm g_e,
 \qquad
 S_e=
 \begin{pmatrix}
 1&0\\
 0&-1\\
 0&1\\
 1&0
 \end{pmatrix},
 \qquad
 P_e=
 \begin{pmatrix}
 1&0&0&0\\
 0&0&1&0
 \end{pmatrix},
 \qquad P_eS_e=I_2.
 }
 \label{eq:SP-massless}
\end{equation}
任何保持关系子空间的原四维算子都可直接约为
\begin{equation}
 O_e^{\rm red}=P_eO_eS_e.
 \label{eq:POS-reduction}
\end{equation}

\subsection{两个顶点作用于同一个二维因子}

约化前，端点 $u_e$ 的常数 h generator 是 $-\ii q_e\sigma_2\otimes I_2$，端点 $v_e$ 的是 $-\ii q_eI_2\otimes\sigma_2$。直接使用式 \eqref{eq:SP-massless} 得到
\begin{align}
 P_e(\sigma_2\otimes I_2)S_e&=\sigma_2,
 \label{eq:first-endpoint-projection}\\
 P_e(I_2\otimes\sigma_2)S_e&=-\sigma_2.
 \label{eq:second-endpoint-projection}
\end{align}
因此四态约为二态后，两个端点不再各有独立 tensor factor，而是在同一个 $\cM_e$ 上分别作用 $+\sigma_2$ 与 $-\sigma_2$。这正是按顶点独立 factorization 失效的唯一 regular 原因。

重要的是，全图 tensor-product 结构没有失效。失效的是分解
\begin{equation}
 \cV_s\stackrel{?}{=}\bigotimes_{v\in V_s}\cV_v,
\end{equation}
因为共享的 $\cM_e$ 不能同时独立归属于 $u_e$ 和 $v_e$；但所有顶点算子仍是同一个全图空间上的 Kronecker 和。

\subsection{contact selector 也随投影直接得到}

论文式 (3.65) 的四态 remaining-term pattern 与 overall scalar 暂时分开后，是
\begin{equation}
 \bm r_e^{(4)}=(0,-1,+1,0)^T.
\end{equation}
它恰好满足
\begin{equation}
 \bm r_e^{(4)}=S_e|1\rangle_e,
 \qquad
 |1\rangle_e=\begin{pmatrix}0\\1\end{pmatrix}.
 \label{eq:contact-vector-projection}
\end{equation}
所以 contact 不需要重新求解：在共享二态中，它就是固定列向量 $|1\rangle_e$。overall scalar $w_e^{(v)}$ 取决于二态 basis。令 $\sigma_e=+1,-1$ 分别表示 $G^{++},G^{--}$；在本节的论文 h-state basis 中，有序第一、第二端点分别是
\begin{equation}
 w_{e,h}^{(u_e)}=+2\ii\sigma_e,
 \qquad
 w_{e,h}^{(v_e)}=-2\ii\sigma_e.
 \label{eq:h-massless-contact-weight}
\end{equation}
第 \ref{sec:package-adapter} 节的常数基变换把它们精确变成 package 的 $-2,+2$。此外只需再乘共同-theta odd-subset 系数与 sector normalization 比值；massless edge 没有额外 Hankel Wronskian prefactor。

\section{含 massless full edge 的全图 regular IBP}

固定 contact-reachable sector $s$。记
\begin{itemize}
  \item $E_0(s)$ 为所有 active 无 theta 一维指数 blocks；
  \item $H(s)$ 为所有仍活动的 massive endpoint h slots；
  \item $E_m(s)$ 为所有未缩并的 massless full edges；
\end{itemize}
约化后的全图空间是
\begin{equation}
 \boxed{
 \cV_s
 =\bigotimes_{b\in E_0(s)}\cE_b
 \otimes\bigotimes_{i\in H(s)}\cH_i
 \otimes\bigotimes_{e\in E_m(s)}\cM_e.
 }
 \label{eq:global-reduced-space}
\end{equation}
因为 $\dim\cE_b=1$，其维数为
\begin{equation}
 \dim\cV_s=2^{|H(s)|+|E_m(s)|}.
\end{equation}

这给出可直接初始化的 building-block registry。每个 block 记录稳定编号、类型、root line/vertex、用户动量名、basis 与 SK/Hankel metadata；每个 root vertex $v$ 记录 incidence list
\begin{equation}
 \cB(v)=
 \{b\in E_0:v(b)=v\}
 \cup\{i\in H:v(i)=v\}
 \cup\{(e,\epsilon_{ve}):e\in E_m,\ v\in e\}.
 \label{eq:vertex-block-incidence}
\end{equation}
一维 blocks 可以在 $M_0$ 中聚合成 $p_{0,v}$，但不能从 registry 删除，因为用户仍按各自的传播子模长和顶点能量输入数值点与微分变量。massless full edge 的同一个编号 $e$ 同时出现在两个顶点的 list 中；这正是“共享而非复制”的数据表达。

为每条有序边 $e=(u_e,v_e)$ 定义 incidence
\begin{equation}
 \epsilon_{ve}=
 \begin{cases}
 +1,&v=u_e,\\
 -1,&v=v_e,\\
 0,&v\notin e.
 \end{cases}
\end{equation}
在论文 h-state basis 中，Hankel branch 选择改变具体解，却不改变一阶系统的矩阵 $-\ii q_e\sigma_2$。因此这里不把 $G^{++}/G^{--}$ 的 branch 符号写进生成元；该符号在换到 package quotient basis 时显式出现。

registry 中每类原子对顶点 $v$ 的贡献可集中写成
\begin{align}
 \Delta M_{1;v}^{(E_b)}&=0,
 &\Delta M_{0;v}^{(E_b)}&=\ii\chi_bq_b I,
 &&b\in E_0(v),
 \label{eq:phase-block-M10}\\
 \Delta M_{1;v}^{(h_i)}&=-(2\nu_i+1)P_1,
 &\Delta M_{0;v}^{(h_i)}&=-\ii k_i\sigma_2,
 &&i\in H_v,
 \label{eq:massive-h-block-M10}\\
 \Delta M_{1;u_e}^{(m_e)}&=\Delta M_{1;v_e}^{(m_e)}=0,
 &\Delta M_{0;u_e}^{(m_e)}&=-\ii q_e\sigma_2,
 &\Delta M_{0;v_e}^{(m_e)}&=+\ii q_e\sigma_2.
 \label{eq:massless-shared-block-M10}
\end{align}
式 \eqref{eq:massive-h-block-M10} 中的 $\nu_i$ 是 2401 的公式参数；本包缺省正 prefactor 统一代入 $\nu_i\to-\nu_i$. 全局 $M_1$ 还要加 component time power $\nu_{0,v}I$ 或 package 的 $A_vI$。式 \eqref{eq:massless-shared-block-M10} 的两个端点作用在同一个二维 slot 上，不能各自再乘出一个新 slot。

在全图空间 \eqref{eq:global-reduced-space} 上，每个顶点的 regular 矩阵直接为
\begin{align}
 \boxed{
 M_{1;v,s}(\nu_{0,v})
 =\nu_{0,v}I_s
 -\sum_{i\in H_v(s)}(2\nu_i+1)\Emb_i(P_1),
 }
 \label{eq:global-M1}\\
 \boxed{
 M_{0;v,s}
 =\ii p_{0,v}I_s
 +\sum_{i\in H_v(s)}\Emb_i(-\ii k_i\sigma_2)
 +\sum_{e\in E_m(s)}\Emb_e(-\ii\epsilon_{ve}q_e\sigma_2).
 }
 \label{eq:global-M0}
\end{align}
massless edge 不进入 $M_1$：缺省取 $\nu=1/2$ 和正 prefactor，把 2401 公式中的参数代为 $-\nu=-1/2$ 后，$2(-\nu)+1=0$。同一个 $\Emb_e(\sigma_2)$ 同时出现在边两端的 $M_{0;u_e,s}$ 和 $M_{0;v_e,s}$ 中，并带相反 incidence 符号。

所有 $M_{0;v,s}$ 都由不同 slot 上的 $\sigma_2$，或同一 shared slot 上的同一个 $\sigma_2$ 构成，所以它们彼此对易，并可由单个全局常数变换
\begin{equation}
 U_s
 =\bigotimes_{i\in H(s)}T_i
 \otimes\bigotimes_{e\in E_m(s)}T_e
 \label{eq:global-T}
\end{equation}
同时对角化；一维指数 slots 只贡献 $[1]$，未显式写出。

这一步不是说任意 Pauli 组合都能这样处理。对任一二维 slot $\alpha$，把各顶点在该 slot 上的局部部分写成
\begin{equation}
 Q_{v\alpha}=c_{0,v\alpha}I+\bm c_{v\alpha}\cdot\bm\sigma.
\end{equation}
存在逐 slot 的共同对角化器所需的条件是
\begin{equation}
 [Q_{v\alpha},Q_{w\alpha}]=0
 \quad\Longleftrightarrow\quad
 \bm c_{v\alpha}\times\bm c_{w\alpha}=0
 \qquad(\text{all }v,w).
 \label{eq:local-common-diagonalization-condition}
\end{equation}
也就是说，同一 slot 上所有 Pauli 方向必须共线。当前论文 h-state 的每个 massive/shared slot 只出现 $\sigma_2$；package quotient 中 massive slots 只出现 $\sigma_2$，每个 massless shared slot 只出现 $\sigma_1$，故全局对角化器改为
\begin{equation}
 U_s^J
 =\bigotimes_{i\in H(s)}T_i
 \otimes\bigotimes_{e\in E_m(s)}H_e,
 \qquad
 H=\frac1{\sqrt2}\begin{pmatrix}1&1\\1&-1\end{pmatrix}.
 \label{eq:package-global-diagonalizer}
\end{equation}
若未来某个 shared slot 同时收到不平行的 $\sigma_1$、$\sigma_2$，单个 $2\times2$ 矩阵也许仍能单独对角化，但不同顶点不再有共同 $U_s$；此时不得继续把大矩阵求逆替换为统一的逐 letter 取倒数。

令 massive bit 为 $a_i\in\{0,1\}$，shared massless-edge bit 为 $r_e\in\{0,1\}$。则
\begin{equation}
 U_sM_{0;v,s}U_s^{-1}
 =\operatorname{diag}\!\left(\ii\kappa_{v,s;\bm a,\bm r}\right),
 \label{eq:global-M0-diagonal}
\end{equation}
其中
\begin{equation}
 \boxed{
 \kappa_{v,s;\bm a,\bm r}
 =p_{0,v}
 +\sum_{i\in H_v(s)}(2a_i-1)k_i
 +\sum_{e\in E_m(s)}
 \epsilon_{ve}(2r_e-1)q_e.
 }
 \label{eq:global-energy-letter}
\end{equation}
同一个 $r_e$ 同时进入 $u_e$ 和 $v_e$ 的 letter，并带相反符号。这是 shared-edge correlation 的完整 regular 内容。

因此 massless quotient 不妨碍对单个顶点使用式 \eqref{eq:global-M1}--\eqref{eq:global-M0}：给定 $v$，只需从 $\cB(v)$ 逐 block 加项。它改变的是不同顶点矩阵之间的支撑关系：$M_{0;u_e,s}$ 与 $M_{0;v_e,s}$ 在同一个 $e$ slot 上分别含相反系数，而不是作用于两个互不相干的 slots。两套逐顶点递推会共享 state bit，却仍各自是闭式矩阵公式，不需要先联立大量 IBP 方程。

\subsection{正/负顶点、SK 符号与 Wick 坐标必须分开}

对每个顶点选定其纯顶点能量指数
\begin{equation}
 \exp(\ii\varsigma_vk_{0,v}\tau_v),
 \qquad \varsigma_v=+1\ \text{或}\ -1,
 \label{eq:vertex-phase-sign}
\end{equation}
并把它作为式 \eqref{eq:phase-exponent-block} 的一个 block。$\varsigma_v$ 是顶点指数符号；它不是 massless propagator 的 $G^{++}/G^{--}$ 标签 $\sigma_e$，也不是有序端点 incidence $\epsilon_{ve}$. 其它无 theta 指数仍按各自 $\chi_bq_b$ 加进 $p_{0,v}$。

正、负顶点对 regular 公式的影响只有
\begin{equation}
 p_{0,v}=\varsigma_vk_{0,v}+\sum_{b\neq b_0(v)}\chi_bq_b,
 \qquad
 M_{0;v,s}\supset\ii p_{0,v}I_s.
 \label{eq:signed-vertex-M0}
\end{equation}
$M_1$ 不含这些一维 blocks；massive h 的 $P_1,\sigma_2$ 原子也不因 Hankel branch 改变。$G^{++}/G^{--}$ 的差别则保存在 theta 分支、Wronskian/contact 权重，以及 package quotient 的 $D_{\sigma_e}$ 中，不能用 $\varsigma_v$ 替代。

现在令 $\tau_v=-t_v<0$. 沿两个顶点各自的阻尼射线统一定义
\begin{equation}
 K_v=\ii\varsigma_vk_{0,v}>0,
 \qquad
 k_{0,v}=-\ii\varsigma_vK_v.
 \label{eq:unified-vertex-wick-coordinate}
\end{equation}
于是
\begin{equation}
 \exp(\ii\varsigma_vk_{0,v}\tau_v)=\exp(-K_vt_v).
\end{equation}
所以 $\varsigma_v=+1$ 时可按参考记号写 $k_{0,v}\to-\ii\,pik_{0,v}$，而 $\varsigma_v=-1$ 时写 $k_{0,v}\to+\ii\,mik_{0,v}$，两者的 $pik_{0,v},mik_{0,v}$ 都是正数。正文只把这些名字用于与文献交流；程序内部应由初始化生成稳定的 $K_v$ 坐标。

用户仍按初始化时自己的 $k_{0,v}$ 和全部传播子模长 $q_e$ 给点。connection 输出到用户坐标时只做可逆 pullback
\begin{equation}
 \dd p_{0,v}=\varsigma_v\dd k_{0,v}+\sum_b\chi_b\dd q_b,
 \qquad
 \dd K_v=\ii\varsigma_v\dd k_{0,v};
 \label{eq:user-to-internal-coordinate-pullback}
\end{equation}
不会把用户的 $q_e$ 擅自改名为 \texttt{pik} 或 \texttt{mik}. contact 把一组 root vertices 合成 $C$ 时，$p_{0,C}=\sum_{v\in C}p_{0,v}$，而阻尼变量满足 $K_C=\sum_{v\in C}K_v$。

\section{直接迭代约化：不先生成大规模 IBP}

\subsection{2401 二进制顺序与递推的精确定义域}

固定 sector $s$，令 $N_s$ 是该 sector 的二维 slots 数，$c_s$ 是 contact 后的 time components 数。沿用 2401 的二进制表示，把 state 写成
\begin{equation}
 \bm a=(a_1,\ldots,a_{N_s}),
 \qquad a_i\in\{0,1\}.
\end{equation}
状态顺序固定为
\begin{equation}
 \bm a(m)=\operatorname{IntegerDigits}(m,2,N_s),
 \qquad 0\leq m<2^{N_s}.
 \label{eq:binary-state-order}
\end{equation}
第一位是最高有效位；Kronecker 积中左侧 slot 慢变、右侧 slot 快变。对任意字串，$\operatorname{Del}_i\bm a$ 表示删除第 $i$ 位，$\operatorname{Ins}_i(\widehat{\bm c},x)$ 表示在第 $i$ 位补入 $x\in\{0,1\}$。这些约定与参考代码的 \texttt{IntegerDigits[...,2,n]}、\texttt{Delete} 和 \texttt{Insert} 完全同序。

MadStree 的一个原生积分可记为
\begin{equation}
 J_s(\bm n;\bm a)
 \equiv\texttt{MSIntegral[}s,\bm n,\bm a\texttt{]},
 \qquad
 \bm n\in\mathbb Z^{c_s},
 \quad
 \bm a\in\{0,1\}^{N_s}.
 \label{eq:MSIntegral-domain}
\end{equation}
这里 $s$ 不是由局部 $\bm n,\bm a$ 反推的标签。设 root propagator 依初始化顺序为
$E=(e_1,\ldots,e_{N_E})$，已收缩的 canonical set 为 $L_s^{\rm ctr}$。定义定长字符串
\begin{equation}
 \operatorname{key}(s)=q_1q_2\cdots q_{N_E},\qquad
 q_j=\begin{cases}0,&e_j\in L_s^{\rm ctr},\\1,&e_j\notin L_s^{\rm ctr}.
 \end{cases}
 \label{eq:sector-bit-string-key}
\end{equation}
不能发生 contact shrink 的传播子恒有 $q_j=1$，top sector 为全 $1$ 字符串。该 key 的类型是字符串，不是二进制整数；前导零不可删除。例如 $E=(e_1,e_2,e_3,e_4)$ 且 $L_s^{\rm ctr}=\{e_1,e_2,e_4\}$ 时，key 是 \texttt{"0010"}。完整积分头始终把该 key 放在第一参数。因此两个 subsectors 即使 $c_s,N_s,\bm n,\bm a$ 全部相同，也仍是不同的
\texttt{MSIntegral} 表达式。同一 contraction set 经不同 contact 路径到达时则合并到同一个 sector。\texttt{MSInitTree} 在返回 context 前生成
\texttt{sectorIdentityCertificate}，一次性检查 canonical contraction sets、sector keys、完整 master expressions、连续 global indices 及完整有序 master 表的 SHA-256 digest；任一碰撞返回
\texttt{SectorIdentityCollision}，不把歧义留给递推或 DE consumer。

因此 \texttt{MSReduce} 的输入域不是“任意 dS 积分”，而是固定初始化 context 中由下列合法对象张成的空间
\begin{equation}
 \boxed{
 \operatorname{span}\!\left\{
 J_s(\bm n;\bm a)\,\middle|\,
 s\in\operatorname{ContactDAG},\
 \bm n\in\mathbb Z^{c_s},\
 \bm a\in\{0,1\}^{N_s}
 \right\}.
 }
 \label{eq:MSReduce-domain}
\end{equation}
也就是说，可以在每个 contact-reachable sector 中任选整数 time shifts 和全部二进制 states，并把有限多个这样的合法对象作线性组合；不能凭这套递推加入新的传播子幂、新的离散指标、context 外 sector 或另一个函数族。

\subsection{包含 contact remaining term 的两向递推}

令 $\bm F_s(\bm n)$ 表示 sector $s$ 的有序 master vector，其中 $n_v$ 平移顶点 $v$ 的论文时间幂次 $\nu_{0,v}$。theta remaining term 写成 lower-sector vectors 的线性组合
\begin{equation}
 \cR_{v,s}(\bm n)
 =\sum_{t\prec s}C_{st}^{(v)}(\bm n)\bm F_t(\bm n_t).
 \label{eq:R-as-CF}
\end{equation}
更具体地，对普通单边 contact 可写成
\begin{equation}
 \cR_{v,s}(\bm n)
 =\sum_{\substack{e\in E_{\rm full}(s)\\v\in e}}
 C_{s,s/e}^{(v,e)}
 \bm F_{s/e}\!\left(\chi_e(\bm n)\right).
 \label{eq:R-edge-sum}
\end{equation}
$\chi_e$ 合并边两端所属的 time shifts，并包含 delta contact 使 child 基准幂少一次的固定平移。这个 parent-to-child 结构对 massive 与 massless edge 完全相同；区别只在局部矩阵 $C^{(v,e)}$。massive edge 使用 2401 式 (3.65) 的互补端点 bits 和符号；massless edge 先作 $4\to2$ quotient，再使用式 \eqref{eq:contact-injection} 的 shared-bit selector 与式 \eqref{eq:package-massless-contact-weight} 的 $-2,+2$。因此 massless 缩并不需要另一套递推公式。

本包的 child base power 特意按 delta contact 固定：$\cR(\bm0)$ 落在 child shift $-1$，而 DE 所需的 $\cR(\bm e_v)$ 恰落在 child shift $0$。所以普通单边 massive/massless contact 的 $R^{(1)}$ 已经直接命中 child master；这里没有“dlog master 选好但迭代约不过去”的问题。

顶点 $v$ 的完整一步 IBP 是
\begin{equation}
 M_{1;v,s}(\nu_{0,v}+n_v)\bm F_s(\bm n-\bm e_v)
 +M_{0;v,s}\bm F_s(\bm n)
 +\cR_{v,s}(\bm n)=0.
 \label{eq:global-one-step-ibp}
\end{equation}

因此降低一步直接为
\begin{align}
 \bm F_s(\bm n-\bm e_v)
 &=A_{-,v,s}(\bm n)\bm F_s(\bm n)
 -M_{1;v,s}^{-1}\cR_{v,s}(\bm n),\\
 A_{-,v,s}(\bm n)
 &=-M_{1;v,s}(\nu_{0,v}+n_v)^{-1}M_{0;v,s}.
 \label{eq:global-Aminus}
\end{align}
升高一步使用对角矩阵
\begin{equation}
 \widetilde M_{0;v,s}=U_sM_{0;v,s}U_s^{-1}
\end{equation}
得到
\begin{align}
 \bm F_s(\bm n+\bm e_v)
 &=A_{+,v,s}(\bm n)\bm F_s(\bm n)
 -U_s^{-1}\widetilde M_{0;v,s}^{-1}U_s
 \cR_{v,s}(\bm n+\bm e_v),\\
 A_{+,v,s}(\bm n)
 &=-U_s^{-1}\widetilde M_{0;v,s}^{-1}U_s
 M_{1;v,s}(\nu_{0,v}+n_v+1).
 \label{eq:global-Aplus}
\end{align}
这就是论文式 (3.37)、(3.47)、(3.50)、(3.66) 在 shared massless-edge space 上的直接推广。需要求逆的仍只有原 basis 中的对角 $M_1$ 和 $U_s$ basis 中的对角 $\widetilde M_0$。现在可直接得到 general-index 的完整奇异层。

在解释奇异层以前，必须把公式中的 $h$ block 与完整物理 mode 区分。若 $d$ 是空间维数，并把曲率耦合等贡献收入 $m_{\rm eff}$，标量指标满足
\begin{equation}
 \nu^2=\frac{d^2}{4}-\frac{m_{\rm eff}^2}{H^2}.
 \label{eq:physical-mode-nu}
\end{equation}
四维 de Sitter 即 $d=3$，完整 mode 含
\begin{equation}
 u_k(\tau)\propto(-\tau)^{3/2}H_\nu^{(1,2)}(-k\tau),
 \label{eq:physical-mode-full-power}
\end{equation}
而不是只有 Hankel 函数。一般 $d$ 维的晚时两支为
\begin{equation}
 u_k(\tau)\sim c_-(-\tau)^{d/2-\nu}
 +c_+(-\tau)^{d/2+\nu}.
 \label{eq:physical-late-time-branches}
\end{equation}
对 minimally coupled massless scalar，$\nu=d/2$，leading branch 恰为零次幂，另一支为 $d$ 次幂。对 $0<m_{\rm eff}^2<d^2H^2/4$ 的轻场，$0<\nu<d/2$，两支实部均为正；对重场 $\nu=\ii\mu$，两支实部都为 $d/2>0$。所以质量项在式 \eqref{eq:physical-mode-nu} 中确实以负号进入，且 generic massive mode 的自然晚时幂次不会自行产生边界发散。

$h=z^{\pm\nu}H_\nu$ 只把式 \eqref{eq:physical-late-time-branches} 的 $\pm\nu$ 幂在“显式 $(-\tau)^A$”与 block 内部之间重新分配；初始化后必须固定同一 convention，但完整物理幂次不变。主积分若保留每条传播子的自然 $(-\tau)^{d/2}$ prefactor，则 massive 两支在 $\tau\to0$ 都衰减。massless leading branch 虽为常数，IBP 所用全微分 primitive 只要另含一个正的显式 $(-\tau)$ 幂便仍趋于零；真正危险的是总 primitive 幂次降到零或负数。这正对应下面 $M_1$ 本征值为零或越过它的层，不能因“massless mode 本身为常数”就额外添加边界项，也不能在零层上继续使用已经丢弃边界项的递推。

对 state $(\bm a,\bm r)$，$M_1$ 本征值是
\begin{equation}
 \lambda_{1;v,\bm a}(n_v)
 =\nu_{0,v}+n_v
 -\sum_{i\in H_v}(2\nu_i+1)a_i,
 \label{eq:M1-general-eigenvalue}
\end{equation}
它与 massless shared bits $\bm r$ 及全部运动学无关。降低 $n_v\to n_v-1$ 时只有
\begin{equation}
 \lambda_{1;v,\bm a}(n_v)=0
 \label{eq:M1-lowering-resonance}
\end{equation}
会产生 $M_1^{-1}$ pole。一般 mixed/Hankel family 仍须结合具体小 $\tau$ 渐近判断；但纯 massless 顶点没有 massive endpoint 项，此时令 $\alpha_v=\nu_{0,v}+n_v$，便有 $M_{1;v}=\alpha_v I$，而相应端点积分正含 $\int_0^\infty\!\dd t_v\,t_v^{\alpha_v-1}e^{-\kappa t_v}=\Gamma(\alpha_v)\kappa^{-\alpha_v}$。因此 $M_1$ 整体在 $\alpha_v=0$ 为零时，未正规化的 Gamma 积分本身也发散，不能直接使用丢弃边界项的 IBP。应先取 $\alpha_v\to\alpha_v+\delta$ 作解析正规化；在合适正规化后的 generic-$\delta$ 系统中 $M_1$ 可逆，递推完成后再展开并取 $\delta\to0$，不会出现物理上的递推失效。对 massive states，若改用尚未吸收 $H\to h$ prefactor 的自然时间幂，则一组 $m$ 个相同阶数、其中 $r$ 个 $a_i=1$ 的贡献为 $r+(m-2r)\nu$；generic 不同阶数不会系统相消，只有 $m$ 为偶数且 $r=m/2$ 时该组 $\nu$ 消去。这类 $M_1$ 零点同样先用时间幂解析 regulator 移开，再在 generic regulator 下递推。

反方向升高 $n_v\to n_v+1$ 时，分母是
\begin{equation}
 \lambda_{0;v,\bm a,\bm r}
 =\ii\kappa_{v,s;\bm a,\bm r},
 \qquad
 \kappa_{v,s;\bm a,\bm r}=0.
 \label{eq:M0-raising-resonance}
\end{equation}
这是运动学 energy-letter 奇异面。这里的 $k_i,q_e$ 都是动量模长，普通矢量动量守恒不迫使 $\kappa$ 为零；当相关动量退化为共线并按轴向分成相反两组 $L,R$ 时，动量守恒化为 $\sum_{\ell\in L}|\bm k_\ell|-\sum_{\ell\in R}|\bm k_\ell|=0$，正是对应 state 符号选择下的 $\kappa=0$。$M_1=0$ 在 raising 公式中只出现在分子，不产生 pole；$M_0=0$ 在 lowering 公式中也只出现在分子。若两者同时为零，则该 row 从两个方向都不能跨越，公式只剩 lower-sector constraint；这时需要另一条关系或另选 basis，但仍不能仅凭 rank loss 宣称一定新增 master。程序应返回方向、shift、state bits、零本征值与未除过的原始 IBP row，不使用固定最大步数掩盖这些层。

0-fold 特例最清楚：$M_1=\nu_0+n$、$M_0=\ii p_0$。lowering 在 $\nu_0+n=0$ 失效，raising 在 $p_0=0$ 失效；这正是论文式 (3.59) 两个分母各自所属的方向，而不是同一个“统一奇异层”。

对单条 massless edge $e$ 缩并到 $t=s/e$，式 \eqref{eq:contact-vector-projection} 直接给出
\begin{equation}
 \boxed{
 C_{s,s/e}^{(v,e)}
 =w_e^{(v)}
 \operatorname{PermuteToGlobalOrder}
 \left(|1\rangle_e\otimes I_{\rm spectators}\right).
 }
 \label{eq:contact-injection}
\end{equation}
若一束 full lines $\mathcal B$ 共享同一个 theta argument，则它们不是独立 delta 的乘积。写
\begin{equation}
 \prod_{e\in\mathcal B}G_e
 =\theta(\Delta)\prod_{e\in\mathcal B}A_e
 +\theta(-\Delta)\prod_{e\in\mathcal B}B_e,
 \qquad
 \mathsf J_e=\frac{A_e+B_e}{2},\quad D_e=A_e-B_e,
\end{equation}
一次求导只给一个 $\delta(\Delta)$，并有
\begin{equation}
 \prod_eA_e-\prod_eB_e
 =\sum_{\substack{\varnothing\neq S\subseteq\mathcal B\\|S|\ {m odd}}}
 2^{1-|S|}\prod_{e\in S}D_e\prod_{e\notin S}\mathsf J_e.
 \label{eq:common-theta-odd-subsets}
\end{equation}
因此一个奇数子集 $S$ 就是一个 simultaneous contact event：所选边同时删除，端点只合并一次，contact selector 是各所选边局部 selector 的张量积再乘 $2^{1-|S|}$。若合并后的未选 massless full edge 两端 coincident，其 odd shared state 为零；未选 massive full edge 使用 $10\sim01$ 的 equal-time canonical quotient。程序先在 raw binary space 构造矩阵，再以 $P_sM^{\rm raw}S_s$ 投影到真实 sector masters，所以 coincident quotient 不要求一般大矩阵求逆。

对 event-reachable component $C$，child 的基准幂可直接由最终 partition 写成
\begin{equation}
 A_{C,s}=\sum_{v\in C}A_v
 +\sum_{e\in E_s(C)}r_e+|C|-1,
 \label{eq:simultaneous-contact-base-power}
\end{equation}
其中 $r_e$ 是被缩并边的 raw contact power。最后一项 $|C|-1$ 是真正发生的顶点合并数；共同 theta 的 $k$ 边 event 只贡献一次，而不是错误的 $k$ 次。式 \eqref{eq:simultaneous-contact-base-power} 只依赖最终 component 与 contracted set，故不同合法 contact 路径得到同一 base power。每个 event 都严格减少 component 数，sector BFS 自动终止，并排除 active self-edge 的再次 contact。

\section{dlog DE 的全图张量公式}

定义每个顶点的对角矩阵值对数
\begin{equation}
 \widetilde\Omega_{0;v,s}
 =\operatorname{diag}_{\bm a,\bm r}
 \left[-\ii\log\kappa_{v,s;\bm a,\bm r}\right].
 \label{eq:global-omega0}
\end{equation}
论文的显式动量项只来自仍为二维 h 系统的 massive endpoint slots：
\begin{equation}
 (\Omega_{{\rm ex},s})_{\bm b\bm a}
 =\delta_{\bm b\bm a}
 \left[-\sum_{i\in H(s)}a_i(2\nu_i+1)\log k_i\right].
 \label{eq:global-omega-ex}
\end{equation}
massless shared-edge slot 不贡献这一项，因为缺省代入后 $2(-\nu)+1=1-2\nu=0$。

固定 sector 的 regular dlog potential 是论文式 (3.55) 的全图版本：
\begin{equation}
 \boxed{
 \Omega_{ss}^{\rm reg}
 =\Omega_{{\rm ex},s}
 -\ii\sum_{v\in V_s}
 U_s^{-1}\widetilde\Omega_{0;v,s}U_s
 M_{1;v,s}(\nu_{0,v}+1).
 }
 \label{eq:global-diagonal-omega}
\end{equation}
它的 letters 就是全部式 \eqref{eq:global-energy-letter}，不需要先对 IBP 关系做线性求解。

因此 dlog producer 与迭代 producer 使用的是同一份 registry，而不是另一套算法：一维 blocks 给 $p_{0,v}$，massive blocks 同时给 $M_1$、$M_0$ 和 $\Omega_{\rm ex}$，massless shared blocks给两个端点的同一 state bit、相反 incidence 和 contact selector。先按式 \eqref{eq:local-common-diagonalization-condition} 构造 $U_s$，再逐 state 写出 $\log\kappa$，最后乘 $M_1(A+1)$，就得到式 \eqref{eq:global-diagonal-omega}。这里从未对 $2^N\times2^N$ 的稠密 $M_0$ 做一般求逆；实际求逆只是对式 \eqref{eq:global-M0-diagonal} 的标量 letters 逐个取倒数。

若改用用户坐标或 Wick 坐标，不重新推导 DE，只对 differential one-form 作式 \eqref{eq:user-to-internal-coordinate-pullback} 的 pullback。若改用 package quotient 或物理导数 basis，则分别使用常数相似变换 \eqref{eq:global-package-similarity} 或含导数项的 gauge transformation \eqref{eq:physical-DE-gauge}；三者不能混成同一个无导数的 basis change。

DE 非对角块必须使用升幂关系中的 remaining term，即论文的 $R^{(1)}$，不能拿 $R^{(0)}$ 的 lower 时间幂次先行约化后替代。本包采用 $(-\tau)^A$，其升幂行是
\begin{equation}
 M_{0;v,s}\bm F_s(\bm e_v)
 =M_{1;v,s}(\nu_{0,v}+1)\bm F_s(0)-\cR_{v,s}(\bm e_v).
 \label{eq:package-raising-contact-sign}
\end{equation}
因此参数微分中的 contact 项为 $+\ii M_0^{-1}\cR$；把论文的 $\tau^\nu$ 写法适配到这里时，不能漏掉这个整体符号。若
\begin{equation}
 \cR_{v,s}(\bm e_v)=\sum_{t\prec s}C_{st}^{(v,1)}\bm F_t(0),
\end{equation}
则论文式 (3.68) 直接推广为
\begin{equation}
 \boxed{
 \Omega_{st}^{\rm contact}
 =+\ii\sum_{v\in V_s}\chi_v(E_{st})
 U_s^{-1}\widetilde\Omega_{0;v,s}U_s
 C_{st}^{(v,1)},
 \qquad t\prec s.
 }
 \label{eq:global-offdiagonal-omega}
\end{equation}
其中 $E_{st}$ 是该 simultaneous event 实际选中的边集，$N_0(E_{st})$ 是其中 masslessFull 的条数，并定义
\begin{equation}
 \chi_v(E_{st})=
 \begin{cases}
  \varsigma_v,&N_0(E_{st})=0,\\
  (-1)^{N_0(E_{st})},&N_0(E_{st})>0.
 \end{cases}
 \label{eq:contact-event-phase}
\end{equation}
pure-massive 的 $\varsigma_v$ 是初始化时固定的顶点指数符号；含 massless 的 event 还包含指数 quotient 相对 h/Wronskian contact 原子的相位。线型相关符号已经全部包含在 $C_{st}^{(v,1)}$：massive 端点使用补位 selector，massless quotient 使用式 \eqref{eq:package-massless-contact-weight} 的 $-2,+2$。单条 pure-massless full edge 的解析 chamber 积分对两个顶点能量求导后，式~\eqref{eq:contact-event-phase} 使完整 DE residual 逐分量精确为零；三条共同 theta 的 simultaneous event 取 $N_0=3$。contact 合并后的 pure-massive component 必须只有一个 $\varsigma$；若输入使同一 component 混合正负指数，公式入口 fail closed。

\subsection{2401 式 (3.68) 的删位、补位与 massless quotient}

先写 2401 的 pure-massive 两顶点结果。设被缩并边在左、右顶点的端点位置分别为 $i,j$；parent states 为 $(\bm a,\bm b)$，child states 为 $(\widehat{\bm c},\widehat{\bm d})$。定义
\begin{equation}
 K_u^{2401}=T_{n_u}^{-1}\widetilde\Omega_{0;u}T_{n_u},
 \qquad
 K_v^{2401}=T_{n_v}^{-1}\widetilde\Omega_{0;v}T_{n_v}.
 \label{eq:component-log-kernel}
\end{equation}
则论文式 (3.68) 的两个端点项逐分量为
\begin{align}
 \left(\Omega_{s,s/e}^{2401}\right)_{(\bm a,\bm b),(\widehat{\bm c},\widehat{\bm d})}
={}&-\ii
 (K_u^{2401})_{\bm a,\operatorname{Ins}_i(\widehat{\bm c},1-b_j)}
 (-1)^{b_j}
 \delta_{\operatorname{Del}_j\bm b,\widehat{\bm d}}
 \notag\\
 &-\ii
 (K_v^{2401})_{\bm b,\operatorname{Ins}_j(\widehat{\bm d},1-a_i)}
 (-1)^{a_i}
 \delta_{\operatorname{Del}_i\bm a,\widehat{\bm c}}.
 \label{eq:2401-binary-top-to-sub}
\end{align}
这就是“删去被缩并端点 bit，再在另一侧补入其互补 bit”的精确含义；所有未删 bits 都由 Kronecker delta 保持。本文的 $(-\tau)^A$ normalized-$J$ convention 把式 \eqref{eq:2401-binary-top-to-sub} 的整体 $-\ii$ 按式 \eqref{eq:package-raising-contact-sign} 改成式 \eqref{eq:global-offdiagonal-omega} 的 $+\ii\chi_v(E)$；2401 的 fixed positive pure-massive 顶点取 $\chi_v=\varsigma_v=+1$。不能只改二进制 selector 而漏掉这个 convention 变换。

现在缩并一条 massless edge $e$。由于 shared slot 使逐顶点 factorization 失效，先在完整 parent sector 上定义全局 kernel
\begin{equation}
 \mathcal K_{v,s}=U_s^{-1}\widetilde\Omega_{0;v,s}U_s.
 \label{eq:global-component-log-kernel}
\end{equation}
在 quotient 后的 parent 全局二进制字串中，该边只有一个 shared bit，记其位置为 $p_e$；child 字串记为 $\widehat{\bm c}$。package basis 的 contact 矩阵分量是
\begin{equation}
 \left(C_{s,s/e}^{(v,e)}\right)_{\bm d,\widehat{\bm c}}
 =\eta_{s,e}\,w_e^{(v)}
 \delta_{d_{p_e},1}
 \delta_{\operatorname{Del}_{p_e}\bm d,\widehat{\bm c}},
 \label{eq:massless-contact-binary-component}
\end{equation}
其中 $\eta_{s,e}$ 收集已吸收的 sector-normalization 与 pinch prefactor，$w_e^{(u_e)}=-2$、$w_e^{(v_e)}=+2$。于是式 \eqref{eq:global-offdiagonal-omega} 中原本的中间态求和直接塌缩：
\begin{equation}
 \sum_{\bm d}
 (\mathcal K_{v,s})_{\bm a,\bm d}
 \delta_{d_{p_e},1}
 \delta_{\operatorname{Del}_{p_e}\bm d,\widehat{\bm c}}
 =
 (\mathcal K_{v,s})_{\bm a,
 \operatorname{Ins}_{p_e}(\widehat{\bm c},1)}.
 \label{eq:massless-binary-sum-collapse}
\end{equation}
因此单条 massless edge 的完整 top-to-sub 块是
\begin{align}
 \boxed{
 \left(\Omega_{s,s/e}^{(e)}\right)_{\bm a,\widehat{\bm c}}
 =-\ii\eta_{s,e}\Big[
 w_e^{(u_e)}(\mathcal K_{u_e,s})_{\bm a,\operatorname{Ins}_{p_e}(\widehat{\bm c},1)}
 +w_e^{(v_e)}(\mathcal K_{v_e,s})_{\bm a,\operatorname{Ins}_{p_e}(\widehat{\bm c},1)}
 \Big]
 }
 \notag\\
 =2\ii\eta_{s,e}\Big[
 (\mathcal K_{u_e,s})_{\bm a,\operatorname{Ins}_{p_e}(\widehat{\bm c},1)}
 -(\mathcal K_{v_e,s})_{\bm a,\operatorname{Ins}_{p_e}(\widehat{\bm c},1)}
 \Big].
 \label{eq:massless-binary-top-to-sub}
\end{align}
所以并不是保留两个独立的 $2\times2$ 端点指标再求和：$4\to2$ quotient 已把它们变成一个 shared bit；所谓“两端点求和”只剩式 \eqref{eq:massless-binary-top-to-sub} 中两个 component log kernels 的相加。

若另一条 massless edge $f\neq e$ 没有被缩并，它的 shared bit 是 spectator，并在式 \eqref{eq:massless-contact-binary-component} 中自动给出
\begin{equation}
 \delta_{r_f^{\rm parent},r_f^{\rm child}}.
 \label{eq:massless-spectator-delta}
\end{equation}
它不产生两个端点 bits 的额外求和；它只继续以相反 incidence 符号进入两个端点的 energy letters 和 $\mathcal K_{v,s}$。mixed massive/massless 情形只需把 massive endpoint bits 与其它 shared bits 按式 \eqref{eq:binary-state-order} 合成同一个全局字串，删除 $p_e$ 后其余 bits 全部由式 \eqref{eq:massless-contact-binary-component} 的第二个 delta 保持。

若 generalized 或 multi-edge contact 使 $R^{(1)}$ 落到 child 的非零 shift $\bm q$，程序现在逐个 child state 调用同一公式递推，构造
\begin{equation}
 \bm F_t(\bm q)=\mathcal Q_t(\bm q\to0)\,\bm F_{\rm all}(0),
 \qquad
 \Omega_{s\to *}^{(v)}
 =-\mathcal K_{v,s}C_{st}^{(v,1)}\mathcal Q_t(\bm q_v\to0).
 \label{eq:shifted-contact-right-composition}
\end{equation}
这里 $\bm F_{\rm all}(0)$ 按 context 的全局主积分顺序排列，因此 $\mathcal Q_t$ 允许继续落到 $t$ 的更低 contact sectors，而非只保留直接 child。程序逐列保存 residual、未消去的 shifted integrals 与方向敏感奇异层。只有全部列闭合，才把式 \eqref{eq:shifted-contact-right-composition} 写入 connection；否则返回
\begin{center}
 \texttt{contactShiftReductionFailed}.
\end{center}
普通单边 contact 仍由基准幂选择直接给 $\bm q=0$。

按剩余 full-edge 数排序 sector 后，完整 connection 是 block triangular：
\begin{equation}
 \dd
 \begin{pmatrix}
 \bm F_{s_0}\\\bm F_{s_1}\\\vdots
 \end{pmatrix}
 =\dd
 \begin{pmatrix}
 \Omega_{s_0s_0}&\Omega_{s_0s_1}&\cdots\\
 0&\Omega_{s_1s_1}&\cdots\\
 \vdots&\vdots&\ddots
 \end{pmatrix}
 \begin{pmatrix}
 \bm F_{s_0}\\\bm F_{s_1}\\\vdots
 \end{pmatrix}.
 \label{eq:block-triangular-de}
\end{equation}
矩阵写成 upper 还是 lower triangular 只取决于 sector 排序；物理内容是 contact 只从 parent 指向严格 lower sector。

\section{无量纲 h 方程、物理二次动量项与实数化}
\label{sec:physical-k2}

式 \eqref{eq:massless-z-eom} 是 building block 内部应使用的方程：
\begin{equation}
 \partial_z^2h_0=-h_0.
\end{equation}
用式 \eqref{eq:tau-z-chain} 恢复物理时间导数才得到
\begin{equation}
 \boxed{
 \partial_\tau^2h_0
 =k^2\partial_z^2h_0
 =-k^2h_0.
 }
 \label{eq:physical-second-derivative}
\end{equation}
所以 $-k^2$ 不是另一套 EOM，而是无量纲常数 $-1$ 经过两次链式法则后的物理系数。内部继续使用 $h_0,h_1$ 二态，就能原样复用论文的 $2\times2$ 公式。

若 backend 实数化采用
\begin{equation}
 k=-\ii\kappa,
\end{equation}
则
\begin{equation}
 -k^2=+\kappa^2.
\end{equation}
若实变量的程序 token 恰好命名为 \texttt{ik}，输出中的 \texttt{ik\^2} 表示这个实变量的平方，不是虚数单位 $\ii$ 乘 $k^2$。

现在区分论文 h-state 与物理端点导数。令
\begin{equation}
 \mathcal F_{ab}
 =\partial_{\tau_u}^{a}\partial_{\tau_v}^{b}G_e
 \quad(a,b=0,1)
\end{equation}
表示尚未把 $-k$ 除掉的物理状态。regular top-sector 部分满足
\begin{equation}
 \mathcal F_{01}=-\mathcal F_{10},
 \qquad
 \mathcal F_{11}=k^2\mathcal F_{00},
 \label{eq:physical-massless-relations}
\end{equation}
而完整 full propagator 的第二式还要加 theta 导数给出的 contact distribution，它正是式 \eqref{eq:contact-injection} 所属的 lower sector。

若选择物理二态 $\bm g_{\rm phys}=(\mathcal F_{00},\mathcal F_{10})^T$，则与论文 h 二态的关系是
\begin{equation}
 \bm g_{\rm phys}=N(k)\bm g_h,
 \qquad
 N(k)=\begin{pmatrix}1&0\\0&-k\end{pmatrix}.
 \label{eq:physical-basis-N}
\end{equation}
物理四态的 regular embedding 因而是
\begin{equation}
 \begin{pmatrix}
 \mathcal F_{00}\\\mathcal F_{01}\\\mathcal F_{10}\\\mathcal F_{11}
 \end{pmatrix}_{\!\rm reg}
 =
 \begin{pmatrix}
 1&0\\
 0&-1\\
 0&1\\
 k^2&0
 \end{pmatrix}
 \bm g_{\rm phys}.
 \label{eq:physical-four-to-two}
\end{equation}
因此若最终 DE 选择物理二态，而不是论文 h 二态，必须做
\begin{equation}
 \boxed{
 \mathcal A_{\rm phys}
 =\dd N\,N^{-1}+N\mathcal A_hN^{-1},
 \qquad
 \dd N\,N^{-1}
 =\begin{pmatrix}0&0\\0&\dd\log k\end{pmatrix}.
 }
 \label{eq:physical-DE-gauge}
\end{equation}
推荐公式 producer 内部始终保留无量纲 h-state basis，使 $S_e,P_e,T,U_s$ 都是运动学无关常数；只有输出物理导数状态时才应用 $N(k)$ 和式 \eqref{eq:physical-DE-gauge}。

\section{与 package convention 的接口}
\label{sec:package-adapter}

\subsection{package 的 massless quotient basis}

前文的 $F_{ab}$ 是论文 h-state。package 为了把每次端点导数的动量因子与离散态分开，使用另一组代数 quotient。对有序边 $e=(u_e,v_e)$，令 $\Delta=\tau_{u_e}-\tau_{v_e}$，并用 $\sigma_e=+1,-1$ 分别表示 $G^{++},G^{--}$，定义
\begin{align}
 \mathsf m_{e,0}&=\theta(\Delta)e^{-\ii\sigma_eq_e\Delta}
       +\theta(-\Delta)e^{+\ii\sigma_eq_e\Delta},\\
 \mathsf m_{e,1}&=-\theta(\Delta)e^{-\ii\sigma_eq_e\Delta}
       +\theta(-\Delta)e^{+\ii\sigma_eq_e\Delta}.
 \label{eq:package-M01}
\end{align}
package 的共享二态与四态 embedding 是
\begin{equation}
 \bm g_e^J=(\mathsf m_{e,0},\mathsf m_{e,1})^T,
 \qquad
 \begin{pmatrix}F^J_{00}\\F^J_{01}\\F^J_{10}\\F^J_{11}\end{pmatrix}
 =S_e^J\bm g_e^J,
 \qquad
 S_e^J=
 \begin{pmatrix}
  1&0\\0&-1\\0&1\\-1&0
 \end{pmatrix}.
 \label{eq:package-S-massless}
\end{equation}
所以 package 中确实是
\begin{equation}
 F^J_{01}=-F^J_{10},
 \qquad
 F^J_{11}=-F^J_{00}.
 \label{eq:package-quotient-relations}
\end{equation}
这个负号不与式 \eqref{eq:h-state-massless-relations} 冲突：$F^J_{ab}$ 已抽出端点导数的 $\pm\ii\sigma_eq_e$，不是 $h_a h_b$。在共同的运动学无关 normalization 下，两种共享二态满足
\begin{equation}
 \bm g_e^J=D_{\sigma_e}\bm g_e^h,
 \qquad
 D_{\sigma_e}=\operatorname{diag}(1,\ii\sigma_e).
 \label{eq:h-to-package-basis}
\end{equation}
因此论文 h-state 中两个端点的 $\mp\ii q_e\sigma_2$，在 package quotient 中正好成为
\begin{align}
 Q^J_{u_e,e}
 &=D_{\sigma_e}(-\ii q_e\sigma_2)D_{\sigma_e}^{-1}
   =+\ii\sigma_eq_e\sigma_1,\\
 Q^J_{v_e,e}
 &=D_{\sigma_e}(+\ii q_e\sigma_2)D_{\sigma_e}^{-1}
   =-\ii\sigma_eq_e\sigma_1,
 \qquad
 \sigma_1=\begin{pmatrix}0&1\\1&0\end{pmatrix}.
 \label{eq:package-endpoint-generators}
\end{align}
等价的完整端点规则是
\begin{align}
 \partial_{\tau_{u_e}}\mathsf m_{e,n}
 &=+\ii\sigma_eq_e\mathsf m_{e,1-n}-2n\,\delta(\Delta),\\
 \partial_{\tau_{v_e}}\mathsf m_{e,n}
 &=-\ii\sigma_eq_e\mathsf m_{e,1-n}+2n\,\delta(\Delta).
 \label{eq:package-massless-time-rules}
\end{align}
这同时给出 package quotient 中有序第一、第二端点的 contact 权重
\begin{equation}
 w_{e,J}^{(u_e)}=-2,
 \qquad
 w_{e,J}^{(v_e)}=+2.
 \label{eq:package-massless-contact-weight}
\end{equation}
它们与式 \eqref{eq:h-massless-contact-weight} 完全一致，因为
\begin{equation}
 D_{\sigma_e}|1\rangle_e\,(+2\ii\sigma_e)=-2|1\rangle_e,
 \qquad
 D_{\sigma_e}|1\rangle_e\,(-2\ii\sigma_e)=+2|1\rangle_e.
 \label{eq:contact-basis-conversion}
\end{equation}
cycle line 把每次 $q_e$ 保存为 $b_e\to b_e-1$；fixed/non-loop line 没有 $b_e$ 槽，必须显式保留结构化模长 $q_e$。两次同端点 regular 导数因而是 $-q_e^2$，两端点各一次则是 $+q_e^2$。

package quotient 的对角化矩阵可直接由论文矩阵得到：
\begin{equation}
 T_e^J=T D_{\sigma_e}^{-1},
 \qquad
 T_e^JQ^J_{u_e,e}(T_e^J)^{-1}
 =T(-\ii q_e\sigma_2)T^{-1}.
 \label{eq:package-diagonalizer}
\end{equation}
所以两套 basis 给出同一组 energy letters；$\sigma_e$ 只交换 shared bit 的标签。若改用固定实 Hadamard 矩阵，这个交换会显式表现为 letter 中的 $\sigma_e(2r_e-1)q_e$。

对整个 sector，把式 \eqref{eq:h-to-package-basis} 嵌入每条未缩并 massless full-edge slot，记所得常数矩阵为
\begin{equation}
 B_s=\bigotimes_{e\in E_m(s)}D_{\sigma_e},
 \label{eq:global-package-basis}
\end{equation}
其余 slots 补单位矩阵。于是全部 regular 与 contact 块可直接变换为
\begin{equation}
 M_{0;v,s}^{J}=B_sM_{0;v,s}^{h}B_s^{-1},
 \qquad
 \Omega_{ss}^{J}=B_s\Omega_{ss}^{h}B_s^{-1},
 \qquad
 C_{st}^{J}=B_sC_{st}^{h}B_t^{-1}.
 \label{eq:global-package-similarity}
\end{equation}
最后一式是左右维数不同的 basis change，不是相似变换；当 $s\to t$ 删除边 $e$ 时，左乘的 $D_{\sigma_e}$ 正好完成式 \eqref{eq:contact-basis-conversion}。

\subsection{初始化时固定 quotient 或冗余 h 表示}

MadStree 对每条 \texttt{masslessFull} 在初始化时固定一次表示。缺省
\texttt{masslessRepresentation -> "Quotient"} 使用一个共享二态
\texttt{masslessShared} slot；显式选择 \texttt{"RedundantH"} 则使用两个
\texttt{masslessEndpointH} slots，按有序端点保留 $(00,01,10,11)$ 四态。
\texttt{"RedundantH"} 只接受 $\nu=1/2$。context 建立后，递推、contact、dlog 与数值层
都固定使用该表示。

冗余表示采用与 quotient 相容的常数 normalization，不额外带入裸 Hankel 乘积的 $2/\pi$。
用式 \eqref{eq:SP-massless} 的 $S_e,P_e$，定义
\begin{equation}
 S_{h\leftarrow J}=S_eD_{\sigma}^{-1},\qquad
 P_{J\leftarrow h}=D_{\sigma}P_e,\qquad
 P_{J\leftarrow h}S_{h\leftarrow J}=I_2.
 \label{eq:madstree-redundant-massless-map}
\end{equation}
对 $G^{++}$，这给出
\begin{equation}
 S_{h\leftarrow J}=\begin{pmatrix}1&0\\0&\ii\\0&-\ii\\1&0\end{pmatrix},
 \qquad
 P_{J\leftarrow h}=\begin{pmatrix}1&0&0&0\\0&0&\ii&0\end{pmatrix}.
\end{equation}
冗余四态的第一个端点 contact 列为
$2\ii\sigma(0,-1,1,0)^T$，第二个端点取反；用
$P_{J\leftarrow h}$ 投影后恰为式 \eqref{eq:package-massless-contact-weight} 的 $-2,+2$。
对含 spectator slots 的 sector，只在对应 massless 两个 endpoint 因子上插入式
\eqref{eq:madstree-redundant-massless-map}，其余因子乘单位矩阵。因此两种 context 都由同一张量公式直接生成
$M_1,M_0$、迭代约化和完整 block-triangular dlog，而不是在验证程序中另写一套 IBP producer。

\subsection{\texorpdfstring{$(-\tau)^A$ 与 normalized-$J$}{(-tau)A 与 normalized-J}}

前文使用论文的 $\tau^{\nu_0}$ convention。package 的唯一公开积分是 normalized object
\begin{equation}
 J_s=\cN_sI_s,
\end{equation}
时间幂写为 $(-\tau_v)^{A_{v,s}}$，所以
\begin{equation}
 \partial_{\tau_v}(-\tau_v)^{A_{v,s}}
 =-A_{v,s}(-\tau_v)^{A_{v,s}-1}.
 \label{eq:package-time-power-derivative}
\end{equation}
在同一个 h-state basis 中，package adapter 因而把论文的式 \eqref{eq:global-one-step-ibp} 精确写成
\begin{equation}
 -M_{1;v,s}(A_v)\bm J_s(\bm a-\bm e_v)
 +M_{0;v,s}\bm J_s(\bm a)
 +\cR^{J}_{v,s}(\bm a)=0,
 \label{eq:package-one-step-ibp}
\end{equation}
其中负号完全来自式 \eqref{eq:package-time-power-derivative}。若使用 package quotient basis，则只需在每条 massless shared slot 上用式 \eqref{eq:h-to-package-basis} 对 $M_0$ 做常数相似变换；$M_1$ 不含该 slot。这个适配不改变 shared-edge slot registry。

因此 package convention 下两向一步公式明确为
\begin{align}
 \bm J_s(\bm n-\bm e_v)
 &=M_{1;v,s}(A_v+n_v)^{-1}
 \left[M_{0;v,s}\bm J_s(\bm n)+\cR^J_{v,s}(\bm n)\right],
 \label{eq:package-lowering}\\
 \bm J_s(\bm n+\bm e_v)
 &=M_{0;v,s}^{-1}
 \left[M_{1;v,s}(A_v+n_v+1)\bm J_s(\bm n)
 -\cR^J_{v,s}(\bm n+\bm e_v)\right].
 \label{eq:package-raising}
\end{align}
这把式 \eqref{eq:M1-lowering-resonance}--\eqref{eq:M0-raising-resonance} 的方向性保留下来；整体符号改变不改变哪一个矩阵在该方向被求逆。

normalized basis 中还必须使用
\begin{equation}
 \partial_xJ_s
 =(\partial_x\log\cN_s)J_s+\cN_s\partial_xI_s,
 \qquad
 C^{J}_{st}=\frac{\cN_s}{\cN_t}C^{I}_{st}.
 \label{eq:normalization-adapter}
\end{equation}
所以 sector 对角块增加 $\dd\log\cN_sI_s$，非对角 contact 块乘 $\cN_s/\cN_t$。只有当该比值已由 basis normalization 变成运动学无关常数时，式 \eqref{eq:global-offdiagonal-omega} 才直接是 constant-residue dlog；否则需要额外 basis/gauge 变换，不能只凭字面上的 $\log\kappa$ 宣称完成。

\subsection{H/h 状态向量的正反变换}

令 $p=s_\nu\nu$。同一个函数 $h$ 与裸 Hankel 状态的局部关系为
\begin{equation}
 \begin{pmatrix}h(\nu,0;z)\\h(\nu,1;z)\end{pmatrix}
 =T_{H\to h}(p,z)
 \begin{pmatrix}H_\nu(z)\\\partial_zH_\nu(z)\end{pmatrix},
 \qquad
 T_{H\to h}(p,z)=
 \begin{pmatrix}
  z^p&0\\p z^{p-1}&z^p
 \end{pmatrix}.
 \label{eq:package-H-to-h}
\end{equation}
其逆矩阵是
\begin{equation}
 T_{h\to H}(p,z)=
 \begin{pmatrix}
  z^{-p}&0\\-p z^{-p-1}&z^{-p}
 \end{pmatrix},
 \qquad
 T_{h\to H}T_{H\to h}=I_2.
 \label{eq:package-h-to-H}
\end{equation}
MadStree 的 \texttt{MSHTohMatrix[nu,z,context]}、\texttt{MShToHMatrix[nu,z,context]} 直接实现这两个矩阵。\texttt{MSConvertBasis} 还可按固定 sector 的 slot registry 把它们作 Kronecker 积：每个 massive endpoint 和 \texttt{masslessEndpointH} 使用自己的 $z=-k\tau_{\rm component}$，massless shared quotient 不是裸 Hankel endpoint pair，故在这一换基中乘单位矩阵。\texttt{NuConvention} 只在 \texttt{MSInitTree} 时选择；context 建立后，局部矩阵与全 sector 换基都固定读取该值，不接受逐调用覆盖。正/负 prefactor context 均已分别检查正反变换为单位矩阵。

这里变换的是被积函数的同序状态向量，不是已经积分后的 \texttt{MSIntegral}。原因是式 \eqref{eq:package-H-to-h} 的 $z^p$ 与 $z^{p-1}$ 会同时改变 component base time power；首版没有建立独立的 H-family sector metadata，因此积分对象重载明确 fail closed，不能把这一步伪装成只翻转 state bit。

dSIBP 020 的 public time-only 表示为 \texttt{J[sectorKey,timeShifts,stateBits]}，与 MadStree 的 \texttt{MSIntegral[sectorKey,timeShifts,stateBits]} 逐槽对应；\texttt{stateBits} 是全 sector registry 中的 $n_i$ 类离散态，不要求按顶点分组。该对应只转换积分身份，不替换 normalization：两边 context 必须对每个 sector 使用同一 $\cN_s$，且求导时都包含式~\eqref{eq:normalization-adapter} 的 $\partial\log\cN_s$。MadStree 的缺省 massless 输入是 $\nu=1/2$ 和正 prefactor，组装 2401 公式时统一代入 $\nu\to-\nu=-1/2$。这一解析等价不替代 line metadata、theta 或 contact normalization。

\subsection{单顶点函数族、局部逆与完整约化接口}

单顶点函数族不需要构造伪图。紧凑输入为
\begin{verbatim}
ctx = MSInitVertexFamily[<|
  "ki" -> {k0,k1,...},
  "nui" -> {A0,nu1,...},
  "hankelBranches" -> {1,2,...}
|>, NuConvention -> "Positive"];
\end{verbatim}
其中 \texttt{First[ki]} 与 \texttt{First[nui]} 分别是纯顶点指数能量和 $(-\tau)$
基准幂；其余位置一一给出 h block 的动量、$\nu=|\nu|$ 和 Hankel branch。显式数据模型使用
\texttt{energy/}\allowbreak\texttt{timePower/}\allowbreak\texttt{hBlocks/}\allowbreak\texttt{exponentialBlocks}。每个纯指数 block 是一维因子，只平移
effective energy；每个 h block 是二维 massive endpoint slot。该 context 可直接调用
\texttt{MSFormulaMatrices}、\texttt{MSMasterIntegrals}、\texttt{MSDLogDE}、
\texttt{MSBoundaryData}，并通过 \texttt{MSEvaluatePath} 完成单阶段数值计算。
\texttt{NuConvention} 只在初始化时选择，
之后所有接口都固定读取 context。

实现中没有一般 $2^p\times2^p$ 符号逆。对 massive endpoint，$M_0$ 的局部非平凡方向是 $\sigma_2$，直接使用论文常数矩阵 $T$ 及其显式逆；对 massless shared edge，局部方向是 $\sigma_1$，使用自逆 Hadamard 矩阵
\begin{equation}
 H_2=\frac1{\sqrt2}
 \begin{pmatrix}1&1\\1&-1\end{pmatrix},
 \qquad H_2^{-1}=H_2.
 \label{eq:package-hadamard-inverse}
\end{equation}
纯指数因子是一维单位矩阵。全 sector 的 $U_s$ 与 $U_s^{-1}$ 只由这些局部矩阵作 Kronecker 积得到。$M_1$ 在 state-bit basis 已经对角，故 \texttt{msM1Inverse} 只计算
\begin{equation}
 \operatorname{diag}(M_1^{-1})_r=\lambda_{1,r}^{-1};
\end{equation}
$M_0$ 则只在共同对角 basis 计算
\begin{equation}
 M_0^{-1}=U_s^{-1}\operatorname{diag}
 \left((\ii\kappa_r)^{-1}\right)U_s.
 \label{eq:package-local-tensor-inverse}
\end{equation}
所以 massless $4\to2$ quotient 后仍完全适配局部对角化求逆。必要条件是同一个 slot 的所有局部算子只含同一 Pauli 方向；未来若出现不共线方向，程序必须拒绝式 \eqref{eq:package-local-tensor-inverse}，不能假装它们可同时对角化。

完整约化入口是
\begin{verbatim}
MSReduce[linearCombination, context, MasterBasis -> Automatic]
\end{verbatim}
它对每个 \texttt{MSIntegral} 依式 \eqref{eq:package-lowering}--\eqref{eq:package-raising} 把各 component shift 递归移到零，并沿 contact DAG 继续约化 lower sectors；重复子问题由 memo 表复用。关键返回字段为
\begin{verbatim}
result, masterBasis, masterRules, coefficientVector,
nonMasterResidual, remainingShiftedIntegrals,
memoizedIntegralCount, singularLayers
\end{verbatim}
其中 master basis、系数向量和 rules 严格同序，result 是显式主积分线性组合。只有 residual 为零且 shifted-integral 列表为空时状态才是 \texttt{reduced}。\texttt{MasterBasis} 可以是全部 context masters 的任意排列，但不能缺项或重复。singular layers 对每一步记录方向、component 和实际被反演的分母：正 shift 向零只报告 $M_0$ 的 energy letters；负 shift 向零只报告 $M_1$ 本征值。这正是式 \eqref{eq:M1-lowering-resonance}--\eqref{eq:M0-raising-resonance} 的程序合同。

\section{任意 time-only 图的大纯指数能量边界与 blow-up}
\label{sec:boundary-blowup}

\subsection{阻尼变量与辅助一维指数}

对每个 root vertex $v$，选择式 \eqref{eq:vertex-phase-sign} 的纯顶点指数，并使用式 \eqref{eq:unified-vertex-wick-coordinate} 的内部正变量 $K_v$. 因为 $\tau_v=-t_v<0$，两种顶点都满足
\begin{equation}
 \exp(\ii\varsigma_vk_{0,v}\tau_v)=\exp(-K_vt_v),
 \qquad K_v>0.
 \label{eq:wick-damping-energy}
\end{equation}
所以 $K_v\to\infty$ 把积分局域到 $\tau_v\to0^-$. 这与文献 \cite{Chen:2024jms} 第 3 节的单顶点 family 边界相同；论文使用的是 $\varsigma_v=+1$、$k_{0,v}\to-\ii\infty$ chart。

若原被积函数在某个顶点没有纯顶点指数，可以加入一个辅助 block $\exp(\ii\varsigma_vk_{0,v}^{\rm aux}\tau_v)$. 它只把
\begin{equation}
 p_{0,v}\longmapsto p_{0,v}+\varsigma_vk_{0,v}^{\rm aux},
 \qquad
 M_{0;v,s}\longmapsto M_{0;v,s}
 +\ii\varsigma_vk_{0,v}^{\rm aux}I_s .
 \label{eq:auxiliary-energy-M0-shift}
\end{equation}
因此在 generic kinematics 下，它不增加 slot 数和主积分数，也不改变 $M_1$。不过 $K_v^{\rm aux}=0$ 是这个 augmented family 的特殊化；是否仍为普通点必须由全 sector letters 决定，不能由“一维 slot 不增加主积分”推出。

\subsection{严格 time 总序、嵌套 blow-up 与 theta 固化}

直接给所有 root times 选择严格总序，并取相反的阻尼能量总序
\begin{equation}
 K_{\sigma(1)}\gg K_{\sigma(2)}\gg\cdots
 \gg K_{\sigma(V)}\gg1.
 \label{eq:strict-energy-rank}
\end{equation}
这一步不要求传播子 incidence graph 是树。本文所说的 time-only 圈图，是允许该 graph 含圈，但所有 line momenta 仍被视为固定参数，只执行各顶点的时间积分；它不包含 loop-momentum integration。严格总序给每一对 times 都定序，其中非相邻顶点的顺序只选择一个 blow-up chart。

定义嵌套坐标
\begin{equation}
 x_j=\frac{K_{\sigma(j+1)}}{K_{\sigma(j)}}
 \quad(1\leq j<V),
 \qquad
 x_V=\frac{1}{K_{\sigma(V)}}.
 \label{eq:nested-blowup-coordinates}
\end{equation}
于是
\begin{equation}
 \frac{1}{K_{\sigma(j)}}=\prod_{r=j}^{V}x_r,
 \qquad x_1,\ldots,x_V\longrightarrow0.
 \label{eq:nested-energy-inverse}
\end{equation}
两顶点时，若 $K_1\gg K_2$，这就是文献第 4.1 节的 $x=K_2/K_1$、$y=1/K_2$。在原逆变量中，$t_1,t_2,t_1+t_2$ 三个 divisor 同时通过二维原点，属于退化交点；式 \eqref{eq:nested-blowup-coordinates} 把它解析为 normal-crossing chart。

指数局域尺度满足
\begin{equation}
 |\tau_{\sigma(1)}|\ll|\tau_{\sigma(2)}|\ll\cdots
 \ll|\tau_{\sigma(V)}|\ll1.
 \label{eq:tau-rank}
\end{equation}
因为所有 $\tau$ 都从负侧趋于零，大 $K$ 对应小 $|\tau|$；所以 $K_u\gg K_v$ 意味着 $\tau_u>\tau_v$，从而
\begin{equation}
 \theta(\tau_u-\tau_v)\longrightarrow1,
 \qquad
 \theta(\tau_v-\tau_u)\longrightarrow0.
 \label{eq:theta-rank-rule}
\end{equation}
对 $G^{++}$ 或 $G^{--}$，必须对传播子展开式中的每一个显式 theta 按其参数顺序使用式 \eqref{eq:theta-rank-rule}；不能仅凭 SK 标签猜一个统一值。
任意 time graph 中的 theta 因而全部变为 $0$ 或 $1$。若某一项要求循环不等式，例如
\begin{equation}
 \theta(\tau_1-\tau_2)\theta(\tau_2-\tau_3)
 \theta(\tau_3-\tau_1),
\end{equation}
则它在每个严格总序 chart 中自动为零，而不是产生新的边界区域。

\subsection{由用户 preferred 普通点选择首选 rank}

用户可以额外给一个 preferred 起点 $P_*$. 该点仍按初始化时的用户变量提交：包括全部 $k_{0,v}$ 和全部传播子模长 $q_e$；package 先用式 \eqref{eq:user-to-internal-coordinate-pullback} 的坐标映射计算
\begin{equation}
 K_v(P_*)=\ii\varsigma_vk_{0,v}(P_*).
\end{equation}
若 $\operatorname{Re}K_v(P_*)>0$，首选 strict chart 取
\begin{equation}
 \operatorname{Re}K_{\sigma(1)}(P_*)>\cdots>
 \operatorname{Re}K_{\sigma(V)}(P_*),
 \label{eq:preferred-rank-from-point}
\end{equation}
相等时依次以 $|K_v(P_*)|$ 和 root vertex order 确定性破缺。注意顺序不能写反：大 $K_v$ 为避免 $e^{-K_vt_v}$ 指数压低，只支持更小的典型 $t_v=|\tau_v|$，所以式 \eqref{eq:preferred-rank-from-point} 对应
\begin{equation}
 |\tau_{\sigma(1)}|\ll\cdots\ll|\tau_{\sigma(V)}|,
 \qquad
 \tau_{\sigma(1)}>\cdots>\tau_{\sigma(V)}.
\end{equation}

这个选择使边界 chart 的尺度层级与 $P_*$ 最接近，通常能缩短第一段 continuation，并可能减少需要绕开的 letter hypersurfaces；它不是“全局最少奇点”的定理。程序仍须对候选 chart 的完整 top/subsector connection 做 normal-crossing 证书，并验证从 chart 内普通 anchor 到 $P_*$ 的路径不穿任何 letter。若某个 $\operatorname{Re}K_v(P_*)\leq0$，该点不能直接定义阻尼 rank；它仍可作为后续普通点，但首选 chart 由用户显式给出或由其它正阻尼 anchor 决定。

\subsection{Contact subsector 继承同一套 root rank}

在 sector $s$ 中，delta contact 把一组 root vertices 合为 component $C$. 其阻尼能量与 rank 定义为
\begin{equation}
 K_C=\sum_{v\in C}K_v,
 \qquad
 r_C=\max_{v\in C}r_v.
 \label{eq:component-energy-rank}
\end{equation}
严格全序保证 $K_C$ 由 $C$ 中唯一的最大 rank root 主导。delta 合并改变 component base time power 和 sector normalization，但不重选一套 subsector rank；remaining theta 连接两个 components 时，只比较二者的主导 root。

令 $v_C$ 是 $C$ 中的主导 root。任一式 \eqref{eq:global-energy-letter} 的边界形式为
\begin{equation}
 \begin{split}
 \kappa_{C,s;\bm a,\bm r}
 &=-\ii K_C+E_{C,s;\bm a,\bm r}\\
 &=-\ii K_{v_C}
 \left[1+\text{higher monomials in }x
 +\frac{\ii E_{C,s;\bm a,\bm r}}{K_{v_C}}\right],
 \end{split}
 \label{eq:boundary-letter-unit}
\end{equation}
其中方括号在 $x=0$ 等于 $1$. 因而每个 contact-reachable sector 的全部 energy letters 都是 coordinate monomial 乘边界非零 unit；这就是严格 rank blow-up 对任意 time graph 线性 letters 的直接解析。

这个结论事实上比 component letter 更一般。固定 sector 的对角块只含式 \eqref{eq:global-omega0} 的 $\log\kappa$. Top-to-sub 公式 \eqref{eq:global-offdiagonal-omega} 只是把 parent 的同一 $\widetilde\Omega_{0;v,s}$ 乘上与运动学无关的 contact map $C_{st}^{(v,1)}$，所以不产生新的 log argument。若 $R^{(1)}$ 没有直接落在 child master 而需要继续递推，新增分母也只来自 child 的 $M_0^{-1}$，仍是 child component 的同型 energy letters。沿严格下降的 contact DAG 归纳，完整 block-triangular DE 中任一含 $K$ 的分母都具有
\begin{equation}
 D_{\alpha,s}(K,q)
 =Q_{\alpha,s}(q)
 +\sum_{v=1}^{V}c_{\alpha,s;v}K_v,
 \qquad
 \operatorname{supp}c_{\alpha,s}\neq\varnothing,
 \label{eq:general-affine-K-divisor}
\end{equation}
其中 $c_{\alpha,s;v}$ 与运动学无关。对普通 component letter，$c_v$ 就是 component 的 indicator；更一般的常数 contact basis 或 SK phase 只会改变非零常数系数。

令
\begin{equation}
 j_\alpha=\min\{j\mid c_{\alpha,s;\sigma(j)}\neq0\}
\end{equation}
是严格能量总序中该分母 support 的第一个位置。代入式 \eqref{eq:nested-blowup-coordinates} 后可严格因子化为
\begin{equation}
 \begin{split}
 D_{\alpha,s}
 =K_{\sigma(j_\alpha)}\Bigg[&c_{\alpha,s;\sigma(j_\alpha)}
 +\sum_{j>j_\alpha}c_{\alpha,s;\sigma(j)}
   \prod_{r=j_\alpha}^{j-1}x_r\\
 &+Q_{\alpha,s}\prod_{r=j_\alpha}^{V}x_r\Bigg].
 \end{split}
 \label{eq:affine-divisor-blowup-factorization}
\end{equation}
方括号在 $x=0$ 等于非零常数 $c_{\alpha,s;\sigma(j_\alpha)}$. 因而所有这类分母的 strict transform 都是 units；其 pullback divisor 只由 $x_{j_\alpha},\ldots,x_V$ 的 coordinate hyperplanes 构成。重复出现的 hyperplane 是同一个不可约 divisor，不增加 normal 数，所以边界附近至多有 $V$ 个彼此独立的 coordinate normals。这证明了：只要 dlog 与 top-to-sub 公式的分母保持式 \eqref{eq:general-affine-K-divisor} 的仿射线性结构，给所有 $\tau_i$ 严格排序所对应的 blow-up 就消除了所有 $K_i\to\infty$ 同时发生造成的简并奇点。证明与原传播子图是树还是 time-only 圈图无关。

程序仍须给出证书，而不能只依赖上述形式检查 top sector。一个简单而充分的算法是：汇总完整 block-triangular connection、normalization gauge 和坐标 Jacobian 的所有不可约 divisor；代入式 \eqref{eq:nested-blowup-coordinates}；抽出所有 $x_j$ 单项式幂；要求余下 strict transform 在 $x=0$ 非零。若仍有 strict transform 消失，则对去重后的局部定义函数 $f_1,\ldots,f_m$ 检查
\begin{equation}
 \nabla f_i\neq0,
 \qquad
 \operatorname{rank}\!\left(\frac{\partial(f_1,\ldots,f_m)}
 {\partial(x_1,\ldots,x_V)}\right)=m\leq V.
 \label{eq:normal-crossing-test}
\end{equation}
不满足式 \eqref{eq:normal-crossing-test} 就仍是退化多变量奇点。对式 \eqref{eq:general-affine-K-divisor} 的 letters，式 \eqref{eq:affine-divisor-blowup-factorization} 已证明所有 strict transforms 都成为 units；程序中的检查是对公式实现及额外 gauge/normalization 的证书，而不是再次猜测这个结论。

公开接口 \texttt{MSBoundaryChartCertificate} 已实现这张证书。它汇总 \texttt{MSDLogDE} 的全部 sector letters、normalization divisors 与 shifted-contact 递推奇异层，返回每个 divisor 的 dominant rank、coordinate monomial、strict transform 和 $x=0$ 值；同时返回用户能量到 $K_v$、再到 $x_j$ 的完整 Jacobian、逐 sector theta 固化表以及 $P_sS_s=I$ quotient 检查。\texttt{RankOrder->All} 枚举全部 root strict charts。\texttt{MSBoundaryData} 在构造无穷远 Frobenius 数据和有限匹配点前强制消费所选 chart 的证书；任一非仿射新分母、消失的 strict transform、零 Jacobian 或未认证 quotient 都返回 \texttt{BoundaryChartNotCertified}，不继续数值输运。

上述定理也给出了自动化范围。完全不含任何 $K_v$ 的奇点已经显式列在 DE letters 中，不是这个边界 producer 的待解决障碍，也不需要 package 去计算奇点上的发散值。缺省路径规划只接受不穿越这些奇点的普通点折线；若直连段命中奇点，程序返回问题点对，用户应显式给出绕行点。用户若主动把路径终点放在这些奇点上，直接使用 package 给出的 DE 自行选择 blow-up。若额外 gauge transformation 引入了不属于式 \eqref{eq:general-affine-K-divisor} 的新分母，则该变换本身不能自动继承上述证明。

每个 sector 的边界条件必须把该 sector 自己当作 top sector：使用 contact 后的 component base power、normalization 和式 \eqref{eq:component-energy-rank}。先计算 bottom sectors，再沿 sector DAG 向上。固定 theta 后的 homogeneous leading term 因子化成 component vertex-family 边界系数的乘积；但 parent 还可收到 lower-sector solution 的 leading contribution。文献第 4.2 节五维例子中，remaining integral 对原 sector 的第二、第三分量已有非零 leading coefficient，正说明只写 homogeneous product 会漏项。通用 producer 应解 block-triangular Frobenius leading system，或等价地保留 theta/delta 的端点展开，递归得到完整同序边界向量。

\subsection{为什么辅助顶点能量不能缺省跑到零}

由式 \eqref{eq:global-M0-diagonal}，
\begin{equation}
 \det M_{0;C,s}\propto
 \prod_{\bm a,\bm r}\kappa_{C,s;\bm a,\bm r}.
 \label{eq:M0-determinant-letters}
\end{equation}
所以 raising 递推式 \eqref{eq:global-Aplus} 在任一 $\kappa=0$ 上失效。最小的反例就是一条 full edge $e=(u,v)$：top-sector letters 为
\begin{equation}
 p_{0,u}\pm q_e,
 \qquad
 p_{0,v}\pm q_e,
 \end{equation}
在 generic $q_e\neq0$ 时把两个辅助能量设为零仍不消失；可是缩边 sector 没有该二维 slot，并必含
\begin{equation}
 \kappa_{\{u,v\}}=p_{0,u}+p_{0,v}.
 \label{eq:contracted-zero-letter}
\end{equation}
因此当这两个 $p_0$ 只由待移除的辅助 blocks 组成时，完整 top/subsector 系统在辅助能量同时为零处是奇异的。更一般地，任一 reachable component 若在目标点既没有提供非零 signed energy 的 active slot，又没有非零总指数能量，就必有零 letter。

对 $M_1=0$ 的 general 结论已经由式 \eqref{eq:M1-general-eigenvalue}--\eqref{eq:M0-raising-resonance} 给出：一般 mixed state 必须另查端点渐近；纯 massless 的例外已在式 \eqref{eq:M1-lowering-resonance} 后说明，此时 $M_1=0$ 正对应未正规化的 $\Gamma(\alpha_v)$ 发散，必须先引入解析 regulator。边界与递推 producer 应消费正规化后的 generic 参数系统，而不把未正规化的零本征值误报为新的主积分或物理递推失效。

若目标点本身落在 $M_0=0$，则必须先在 generic 顶点能量约化再取受控极限，或由用户基于已输出 DE 另行处理该奇点，不能把零点交给普通点数值输运。

因此推荐流程是：先按式 \eqref{eq:preferred-rank-from-point} 选择首选 chart，从式 \eqref{eq:nested-blowup-coordinates} 的边界得到同序向量，在 chart 内取一个小而有限的普通点，再由用户按所需同伦类给出到 preferred 普通点和最终目标点的多变量折线。缺省规划拒绝任何穿过 letter hypersurface 的直连段，并报告问题点对；用户可增加明确绕行点，只有显式选择奇点折跃时才用局部基穿越奇点并自行确认多值分支。辅助顶点能量若要移除，应逐个取零，并在每一步对所有 sectors 重做普通点证书。只有证书通过时，“先把这个辅助能量跑到零再跑其它参数”才是合法的简化；一旦零点是奇点，自动流程停止并返回该点的 DE letters，由用户自行决定后续 blow-up。

\subsection{当前 2411 无穷远边界与 FlintNDE 输运}

生产入口不在有限 Euclidean 普通点直接计算定义积分。\texttt{MSBoundaryData} 以文献 \cite{Chen:2024jms} 的 $K_v\to\infty$ 数据为唯一边界来源。包括单顶点 massive-external family 在内，所有已闭合 tree/time-only context 都统一消费 sector DAG、component、slot registry、normalization、state order 与 strict rank，并走“边界领头项 + FlintNDE Frobenius 递推”路线。式 (3.44)--(3.46) 的单顶点显式多变量级数只作为测试 oracle，不进入生产分派。

具体地，对 sector $s$ 的 component $C$ 定义
\begin{align}
 \alpha_{sC}(\bm a)&=A_{sC}+1-
 \sum_{r\in H(s,C)}a_r(2\nu_r^{\rm formula}+1),
 \notag\\
 \ell_{sC}&=V-\min_{v\in C}\operatorname{rank}(v)+1,
 \qquad
 K_{sC}(t)=c_{sC}t^{-\ell_{sC}}+O(t^{-\ell_{sC}+1}).
 \label{eq:generic-sector-boundary-component}
\end{align}
这里 $H(s,C)$ 只含该 component 上仍活跃的 Hankel endpoint slots；massless shared slots 不进入 $\alpha_{sC}$，而是按固定 theta chamber 给出 $\pm1$ 的 state factor。对应 master 的 Frobenius 指数和物理 leading weight 为
\begin{align}
 \lambda_{s,\bm a}&=\sum_C\ell_{sC}\alpha_{sC}(\bm a),
 \notag\\
 W_{s,\bm a}&=N_s\Phi_s\Xi_{s,\bm a}
 \prod_C\left[
 (-\ii)^{A_{sC}+1}\Gamma(\alpha_{sC})
 E_{sC,\bm a}\,c_{sC}^{-\alpha_{sC}}
 \right].
 \label{eq:generic-sector-boundary-weight}
\end{align}
$N_s$ 是 sector normalization，$\Phi_s$ 是已初始化的 massive-contact phase，$\Xi_{s,\bm a}$ 收集 shared-massless chamber sign 与 RedundantH normalization，$E_{sC,\bm a}$ 是该 component 全部 Hankel endpoint coefficients 的乘积。所有量都从当前 sector metadata 取得，不能从 root sector 猜 child 的幂次或 slots。

把完整 dlog potential 沿 $K_{\sigma(j)}=B^{V-j+1}t^{-(V-j+1)}$ 拉回并记 residue 为 $R_0$。对每个 $(s,\bm a)$，程序固定该 sector 内对应 root master 的分量为 $1$，非 ancestor sectors 为 $0$，解
\begin{equation}
 (\lambda_{s,\bm a}I-R_0)v_{s,\bm a}=0.
 \label{eq:generic-sector-indicial-vector}
\end{equation}
于是 $W_{s,\bm a}v_{s,\bm a}$ 自动包含 child solution 注入 parent 的 leading 分量。重根或整数差共振不另造图专用参数；完整 exact residue 与这些 $\{\lambda,0,v\}$ 交给 FlintNDE 的 power-log 递推。有限定义积分只作为独立验证 oracle，不进入生产输入。

对两顶点等时间幂 $\nu_0$ 的负 prefactor convention，child master 直接采用 2411 式 (4.2)
\begin{equation}
 I_R=-\frac{4\ii}{\pi}e^{\pi\operatorname{Im}\nu_1}k_s^{-2\nu_1-1}
 \int_{-\infty}^{0}\dd\tau\,(-\tau)^{2\nu_0-2\nu_1}
 e^{\ii(k_{12}+k_{34})\tau}.
 \label{eq:paper-two-vertex-child-normalization}
\end{equation}
其第五分支的 leading vector 与物理系数分别为
\begin{align}
 v_5&=\left(0,\frac{1}{2\nu_1-\nu_0},\frac{1}{\nu_0+1},0,1\right)^T,
 \notag\\
 C_5&=-\frac{4\ii}{\pi}e^{\pi\operatorname{Im}\nu_1}
 k_s^{-2\nu_1-1}e^{\ii\pi(\nu_1-\nu_0)}
 \Gamma(2\nu_0-2\nu_1+1).
 \label{eq:paper-two-vertex-contact-boundary}
\end{align}
正 prefactor context 只把这些公式中的论文参数统一换成 \texttt{formulaNu}$=-|\nu|$。上述五维结果现在只是式 \eqref{eq:generic-sector-boundary-component}--\eqref{eq:generic-sector-indicial-vector} 的一个检查例，不是生产分派；一般不等 time powers 直接读取各 component 自身的 $A_{sC}$。实正物理动量的 massive pinch normalization 固定为式 \eqref{eq:paper-two-vertex-child-normalization} 的 $-4\ii k_s^{-2\nu_1-1}/\pi$，不能把形式符号 $(-k_s)^{-2\nu_1-1}$ 交给 Mathematica 主支。

\texttt{MSEvaluatePath} 在 MadStree 侧按输入顺序把用户点划分为最大连续复仿射单变量段
$x(s)=x_0+s v$，每段只拉回一次单变量矩阵 DE，并把全部 exact 复参数点、边界和选项放入
同一个 UTF-8 请求。MadStree 不选择节点或绕行；FlintNDE 在同一进程内完成边界初始化、
可选路径规划和输运。开启规划时，同一展开节点覆盖的点组成 evaluation bucket：大桶通过
子积树/余数树快速多点求值，小桶使用 iterative 算法；关闭规划时用户点严格作为节点。
不同段不共享局部系数，因而没有构造多变量高维 Taylor 球。后端位置由相对项目根目录的单一配置
变量控制；运行文件与 Python cache 写入调用脚本目录的 \texttt{results\_temp/}，不写 package
源码目录。

任何经过中途节点的多点输运都称为折跃；只有显式穿过奇点并用局部基连接两侧匹配点时才称为
奇点折跃。\texttt{SingularityMode -> "Avoid"} 是缺省模式：相邻用户点连线命中奇点时返回
\texttt{"Singular Path Pair"} 与问题点对；显式选择
\texttt{SingularityMode -> "SingularityJump"} 才允许奇点折跃，并提示用户确认等价绕行类的
多值分支。\texttt{MessageLanguage -> "EN"} 缺省给出英文提示，设为 \texttt{"CN"} 时给出中文
提示；两个值严格区分大小写。每段的有限奇点由后端从拉回后的单变量系统内部识别，用户无需提取
或输入奇点列表。

若某段所有 dlog letters 都是常量，则拉回导数全部为零，该段是无有限极点的零连接并正常输运。
\texttt{MSBoundaryData} 与 \texttt{MSEvaluatePath} 的缺省工作精度为 200 位；正规化自动规划
取 200 位与自适应估计的较大者，用户显式指定精度时直接覆盖。FlintNDE 工作位数为
$\lceil p_{\rm dec}\log_2 10\rceil+32$：缺省 200 位对应 697 bit，70、100 位分别为 265、365 bit，
更高精度继续按同一公式增长；奇点局部基门禁也按当前工作精度收紧。线段投影、匹配比例、旋转
因子和 winding/monodromy 均保持当前 Arb 精度。

\subsection{正规化参数的符号 Laurent 支撑证书与误差控制}

共同正规化参数记为 $\epsilon$。用户调用
\texttt{MSReconstructEpSeries[..., MaximumEpPower -> $m$]} 时只指定需要返回的最高幂 $m$；
最低幂 $n_{\min}$ 不是数值拟合参数。任何非零 $\epsilon$ 的 NDE 任务启动前，MadStree 都从
同一个生产 context 的物理边界与 dlog DE 自动构造支撑证书。具体门禁为：
\begin{enumerate}[label=\arabic*.]
  \item 路径坐标不得依赖 $\epsilon$，且所有 dlog 矩阵、拉回 connection 与 residue 在
  $\epsilon=0$ 不含负 Laurent 幂；
  \item 定义积分给出的每个物理边界分支必须是有限 Laurent 型；具有相同极限 Frobenius
  指数和 log 次数的分支先按物理权重与 leading vector 合并，再逐分量取 valuation；
  \item 合并边界中至少一个最低阶系数必须可证明非零。所有分量都在用户检查范围以上、
  非整数幂、非 Laurent 结构或无法判零时均 fail closed。
\end{enumerate}
若 $n_{\min}<0$，证书表示边界是亚纯的有限 Laurent 型而不是在 $\epsilon=0$ 解析；只有
$n_{\min}\geq0$ 才同时记录解析性。沿当前有限路径，$\epsilon$-解析且可逆的基本解保持非零
边界 formal data 的 valuation，因此该 $n_{\min}$ 可直接用于后续数值采样规划。单个
Frobenius 递推算子在 $\epsilon=0$ 的降秩只作为共振诊断保存；它本身不能越过已经合并的
物理边界而被解释成 $1/\epsilon$ pole。

Example 06 的无质量三顶点采用
$a_1=a_2=a_3=1+\epsilon$。其物理边界权重包含
$\Gamma(2+\epsilon)^3$、$\Gamma(2+\epsilon)\Gamma(4+2\epsilon)$ 和
$\Gamma(6+3\epsilon)$，都在 $\epsilon=0$ 正则；实际九维 DE 也无负 $\epsilon$ 阶，故程序在
任何 NDE 求解前认证 $n_{\min}=0$。这意味着该例没有 pole；程序不会为了演示而虚构
$1/\epsilon$ 项。

认证支撑后，内部方阵插值缺省至少拟合到用户最高阶以上两阶。独立验证点从不进入拟合矩阵；
若其相对残差高于固定的 \texttt{EpGoalDigits} 门槛，下一轮缺省再增加两阶，只求解新增生产点，
并复用所有既有生产点和独立验证点。最多三轮仍未通过时返回失败，不降低容差、不改用最小二乘，
也不从头重算缓存。不同 $\epsilon$ 任务由 \texttt{ParallelTaskCount -> 12} 的缺省外层进程池
调度；实际并发取任务数与上限的较小者。

用户点只接受裸坐标或 \texttt{\{coordinate,}\allowbreak\texttt{"tmp"\}}；前者缺省保存，后者是临时途经点。
普通保存点在 \texttt{MSEvaluatePath} 返回值中携带坐标、完整主积分向量、状态和用户序号；使用
\texttt{MSExportEvaluationData[result,}\allowbreak\texttt{ MSOutputDirectory -> outputDirectory]} 导出 CSV/JSON。
临时点不进入普通点导出。

FlintNDE 在启动前分类局部奇点并执行 capability preflight；MadStree 当前 generic boundary producer 只在拉回 connection 是 exact regular singular 时交付 $\{a,b,C\}$，因此新增 $\{\phi,a,b,C\}$ 是后端路径保存能力而不是改变 dS 边界来源。若输入需要超出后端已声明能力的高阶或不规则奇点处理，适配器直接 fail closed，不把普通 Taylor 输运伪装成奇点解。

原被积函数的 $e^{-K_vt_v}$ 阻尼在 $K_v\to\infty$ 的原坐标表述中看起来具有指数渐近；这不等于正规化主积分的 dlog 系统在 blow-up 坐标 $x_j=0$ 必有本性奇点。抽出式 \eqref{eq:generic-sector-boundary-weight} 的已知物理权重并作 nested blow-up 后，当前接受的 exact 系统必须是普通匹配问题或 Fuchsian regular-singular 问题。若实际拉回系统仍有不可约二阶以上 pole、需要 Stokes/ramification 或不属于当前 exact field，capability preflight 会返回实际分类并停止；package 不把它自动解释成已解决的边界条件。

当前已经端到端检查单顶点三 massive 的 2411 边界、两顶点 pure-massless 的解析 chamber 边界、三顶点 massive/massless 的通用 15 维边界与输运，以及两顶点 $G^{++}$ 五维系统。pure-massless 定义积分同时检查 contact 原子权重的两张 DE residual 和完整目标向量；massive 五维系统的 connection、奇点拉回、五个 leading vectors、五个系数和目标向量均与独立论文路线一致。结构检查还覆盖共同-theta simultaneous events、massless/massive coincident quotient，以及 triangle time-only 圈图的通用 boundary/residue。FlintNDE 适配器另检查普通/正则奇点保存点和超能力 pole 的 fail-closed。代表性运行证据不等于对所有 topology 与参数区域逐点认证；通用性来自同一 metadata 数据流与结构不同拓扑的共同回归，而不是列举专用生产分支。

\section{直接公式 producer 的最小算法与验收}

不运行 sampled IBP 或 reduction 时，公式 producer 只需执行以下步骤：
\begin{enumerate}[label=\arabic*.]
  \item 从 topology 建立全图 building-block registry：每个无 theta 指数一个有独立编号的一维 slot，每个 massive endpoint 一个二维 h slot，每条未缩并 massless full edge 一个共享二维 slot；同时保存每个 vertex 的 block incidence list。
  \item 对 massless full edge 固定式 \eqref{eq:SP-massless} 的 $S_e,P_e$ 和端点 incidence；cross/external exponent 不建立共享 slot。
  \item 按式 \eqref{eq:global-M1}--\eqref{eq:global-M0} 逐顶点组装 Kronecker 和；一维指数项与所有二维项使用同一个 embedding API，并严格区分 $\varsigma_v$、$\sigma_e$ 与 $\epsilon_{ve}$。
  \item 先验证式 \eqref{eq:local-common-diagonalization-condition}，再用一个全局 $U_s$ 列举式 \eqref{eq:global-energy-letter} 的对角元；只反演 $M_1$ 和 $\widetilde M_0$ 的标量对角元。
  \item 按式 \eqref{eq:common-theta-odd-subsets} 对 current-component bundles 做 event BFS；对 coincident spectators 构造 $S_s,P_s$ quotient，并用式 \eqref{eq:simultaneous-contact-base-power} 固定 child powers。
  \item 用式 \eqref{eq:global-Aminus}--\eqref{eq:global-Aplus} 输出 general-index 单步递推；用事件级 contact selector 追加 lower-sector source，并按方向返回 $M_1/M_0$ 奇异层。
  \item 用式 \eqref{eq:global-diagonal-omega}、\eqref{eq:global-offdiagonal-omega} 直接输出 block-triangular dlog candidate；非零 child shift 按式 \eqref{eq:shifted-contact-right-composition} 自动约回全局 masters，再应用式 \eqref{eq:normalization-adapter}。
  \item 对所选或全部 root rank 调用 \texttt{MSBoundaryChartCertificate}；只有完整 connection、normalization、Jacobian、theta fixing 与 quotient 同时通过才进入自动边界。
\end{enumerate}

最低 exact 验收应包括：
\begin{itemize}
  \item 单条 massless full edge 的 $4\to2$ 秩、$P_eS_e=I_2$、式 \eqref{eq:first-endpoint-projection}--\eqref{eq:second-endpoint-projection} 和 contact vector；
  \item 两顶点单边图中，同一个 shared bit 在两个 $\kappa_v$ 中符号相反；
  \item 三顶点 chain 中两个 shared slots 的四个共同 eigenstates，而不是每个顶点独立枚举；
  \item mixed massive/massless tree 的全局 $U_sM_{0;v,s}U_s^{-1}$ 全部对角；
  \item 公式单步关系与 naive time-IBP 在完全相同 ordered normalized basis 下逐项相等；
  \item 公式 DE 与 naive DE 逐矩阵元相等，所有 lower-sector residual 为空，并通过 exact flatness；
  \item mixed 三顶点中，dSIBP 020 的公开 \texttt{ds} 经逐槽 adapter 和 \texttt{MSReduce} 得到的五个 $15\times15$ 矩阵，与直接公式 DE 严格相等；
  \item 物理输出 basis 额外包含式 \eqref{eq:physical-DE-gauge}，没有遗漏 $k$、$k^2$ 或 normalization derivative。
\end{itemize}

\section{最终判断}

纯 massive 公式从一开始就是 building-block registry 上的 Kronecker 和：每个无 theta 指数给 $\Delta M_1=0$、$\Delta M_0=\ii\chi_bq_bI$，每个二维 h endpoint 给 $\Delta M_1=-(2\nu_i+1)P_1$、$\Delta M_0=-\ii k_i\sigma_2$，所有项都乘其它 slots 的单位矩阵。MadStree 缺省正 prefactor convention 统一代入 $\nu_i\to-\nu_i$，所以第一项改为 $-(1-2\nu_i)P_1$。正、负顶点只改变纯顶点指数的 $\chi_b=\varsigma_v$；沿 $k_{0,v}=-\ii\varsigma_vK_v$、$K_v>0$，两者都成为 $e^{-K_vt_v}$，不会改变二维 h 原子。

加入含 theta 的 massless full edge 后，应先保留与 massive 相同的两个端点二态，再用 $S_e,P_e$ 把四态约为一个共享二态。两个顶点不再拥有彼此独立的该边因子，而是在同一个 slot 上以相反符号作用；因此逐顶点 tensor factorization 失效，但全图 Kronecker 和、共同对角化、迭代矩阵和 dlog DE 都仍可直接写出。theta 导数只增加由 $|1\rangle_e$ 生成的 lower-sector block。论文 h-state 的 $F_{11}=+F_{00}$ 与 package quotient 的 $F_{11}^J=-F_{00}^J$ 由常数 $D_{\sigma_e}$ 严格连接，不能混写。

单顶点递推不受共享槽阻碍：从该顶点的关联表逐块相加即可得到 $M_{1;v,s},M_{0;v,s}$。不同顶点只是在同一无质量槽上共享状态位并取相反的端点符号。当前每个二维槽中只出现一个固定泡利方向，故同槽算子两两对易；局部 $2\times2$ 对角化器的 Kronecker 积可把大矩阵求逆化成逐能量字母取倒数。未来若同一槽出现不共线的泡利方向，则必须停用这条快速路径。lowering 只在 $M_1$ 本征值为零时失效，raising 只在 $M_0$ 的能量字母为零时失效；未被该方向求逆的矩阵即使为零，也不产生该方向的 pole。

dlog DE 由同一 registry 直接产生：sector diagonal block 使用共同对角化后的 $\log\kappa$ 与 $M_1(A+1)$，top-to-sub block 使用同一 $\log\kappa$ 乘 $R^{(1)}$ contact map，再应用 sector normalization。因而无需先生成大量 IBP 再解出 DE；但只有 normalization ratio、constant residue、block closure 与 flatness 全部通过后，输出才可认证为 dlog，而不是仅称为 connection candidate。

边界同样使用初始化后的 $K_v=\ii\varsigma_vk_{0,v}$，而不把用户变量永久改名为 \texttt{pik0}/\texttt{mik0}。preferred 普通点按 $\operatorname{Re}K_v$ 降序选择首选 strict chart；这等价于 $|\tau_v|$ 升序，因为大阻尼能量只支持更小的典型 $|\tau_v|$。该选择用于缩短首段 continuation、减少可能需要绕开的 letters，不替代普通点和 normal-crossing 证书。

massless building block 应始终写成无量纲 h 方程 $h_{zz}+h=0$。物理式 $\partial_\tau^2h=-k^2h$ 完全等价，只是 $\partial_\tau=-k\partial_z$ 的两次链式法则；它不要求引入另一套不兼容的局部 EOM。这样既保留论文公式的 $2\times2$ 结构，也能在最终物理 basis 中正确恢复 $k^2$ 和 contact。

\bibliographystyle{unsrtnat}
\bibliography{references}

\end{document}
